% content.tex — the body of paper_02_ai_mediated_intimacy, as a book chapter.
% GENERATED by split_content.py. The standalone paper still builds from
% paper_02_ai_mediated_intimacy.tex; both read this file, so they cannot drift apart.

\begin{abstract}
\noindent
Suppose an AI system, drawing on the recorded history of an intimate relationship, suggests to one partner a small tender gesture: a word of appreciation, a remembered flower. Suppose further that the suggestion works, that the relationship genuinely improves by the joint standard of measurable relational quality and sincere first-person experience, and that no one is deceived at the point of delivery. Is anything wrong? A predecessor paper answered: not a clean acquittal, and suspended the question of why. This paper takes the suspended question as its object. It proceeds dialectically. From Hegel it takes the insight that mediation as such is innocent: spirit, and love, exist only through externalization. From Marx it takes the antithesis: alienation names not a defective product but a process separated from its agent, and the tender tip realizes an ironic form of the 1844 manuscripts' fourth alienation, an estrangement of person from person enacted in the very name of bringing persons closer. From the universalist grammar of modern law it takes the strongest available thesis: instruments do not break causation, agency law exists precisely to let imputation pass through intermediaries, and a reductio (the search engine, the friend's advice, the relationship book) destroys any categorical claim that mediation as such denatures love. From Buddhist dependent origination it takes a still deeper version of the thesis, that no unmediated act ever existed, together with Buddhism's own non-delegable core: practice, attention, the raft that must not be clung to. The synthesis is doctrinal in form: imputation through instruments holds as the default rule, but there exists a criterion-governed exception domain, modeled on the legal category of strictly personal acts (h\"ochstpers\"onliche Handlungen), within which the fiction of imputation fails, not because the tool interrupts causation but because acts of this class do not admit substitutable links at all. Four criteria bound the domain (object of mediation, provenance of data, direction of initiative, attribution expectations), and a fifth, dynamic criterion, convergently derivable from Hegel's externalization-and-return and from the Buddhist simile of the raft, adjudicates the residual cases. Intimate acts are then analyzed as layered phenomena, bearing an allocative face, on which attention is genuinely scarce and budgetable and every allocative apparatus (legal imputation, Robbinsian economics, the labor theory of value) holds full jurisdiction, and a constitutive face, on which what matters is who expends the attention and the same apparatuses suffer category mismatch; the specific danger of AI mediation is named distributive reduction, the optimization of an act's allocative face that silently substitutes its constitutive face. A political economy of the cognitive budget makes the danger precise in allocative grammar's own voice: the machine simultaneously reallocates the budget (legitimately), substitutes the expender of constitutive items (self-defeatingly), migrates the budget's managerial function (a Taylorism of love), and, at the systemic level, deflates the value-form of the tender gesture by collapsing the socially necessary attention time that made gestures credible signals of personal attending. An eliminative experiment then removes capital from the scenario altogether: across commercial, single-built, and jointly built configurations, exploitation and domination vary, while the substitution of attending does not, establishing a wrong insensitive to ownership. The eliminative result is then given its constructive complement in the theory of generative justice: the paper's remedies, self-hosting, attribution hygiene, scaffolded return, and co-constitution, are unified as a single four-layer practice of unalienated value circulation, in which the recipient is recognized as the uncredited co-generator of the gestures directed at her, and which is bounded exactly where the value at issue admits no proxy generation. The single-builder residue is developed into a theory of the architectural constitutionalization of intimacy, a mediating layer as a constitution drafted by one party, exercising third-dimensional power in the form of the good, against which subsequent consent is habitation rather than authorship, and whose only symmetric form is constitution practiced as a verb. A concluding feminist hearing settles the paper's debts to the tradition it has borrowed from throughout: it shows the cognitive budget and the constituent power to be gendered in fact, distinguishes the redistribution of attending from its counterfeit, automated substitution, and supplies, from the ethics of care, the affirmative account on which the attending is not the means of love but its mode of existence. The paper distinguishes two wrongs with two remedies: misattribution, curable by transparency, and process alienation, not curable by consent, and answers the predecessor paper's open question in the negative and qualified form that the dialectical method requires.
\end{abstract}

\section{Introduction: Cyrano's Question}
\label{p02:sec:intro}
% =====================================================================

In the third act of Rostand's play, Roxane stands at her balcony and falls in love with words spoken from the shadows beneath it \parencite{rostand1897cyrano}. The voice is Christian's; the sentences are Cyrano's. Every proposition Roxane hears is, in a narrow sense, true: she is loved, extravagantly, by the totality of what addresses her. And yet the scene is the canonical European image of something gone wrong in the supply chain of devotion. Roxane's error is not factual but attributive. She assigns the authorship of the attending to the wrong node in the chain, and on that misassignment she builds a marriage. The play does not resolve whether she was wronged by lies (there were, strictly, few) or by an arrangement of mediation that made correct attribution impossible. That unresolved question, transposed from verse to jurisprudence, is the subject of this paper.

The transposition is occasioned by a contemporary structure that repeats Cyrano's geometry at scale. A predecessor paper \parencite{huang2025toward} examined, as one module of a family-and-relationship management application, what it called the tender tip: an AI conversational layer which, with reference to relationship history and the popular five-love-languages heuristic \parencite[cited there with explicit qualifications]{chapman1992five}, generates small suggestions intended to nurture the relationship, a prompt to express appreciation, a reminder of a remembered preference. The predecessor paper's verdict was deliberately incomplete. At the point of delivery there is no deception: the user knows the suggestion comes from a machine. But the wrong, if there is one, re-enters downstream. The suggestion is executed as a gesture; the receiving partner attributes the thoughtfulness to the partner who performed it; and the cognitive labor behind the gesture was, in part, the machine's. The predecessor paper named this structure \term{a supply chain of thoughtfulness whose origin the final recipient misreads}, allowed the module provisional survival under four design constraints, and explicitly kept open the question whether those constraints suffice, and more fundamentally, whether a genuinely improved relation can redeem an alienated process.\footnote{\textcite[\S 6.6, \S 6.7]{huang2025toward}. The four constraints: opt-in and infrequent; suggesting occasion and register but not content; never presented as knowledge of the partner's mind; refusable and unrecorded. The concept of \emph{proxy epistemic injustice}, on which this paper builds, is introduced there at \S 6.7 as an extension of \textcite{fricker2007epistemic}.}

This paper takes up the open question in its sharpest form. Stipulate everything in favor of the practice. The tip works. The relationship improves, and improves by the demanding joint standard the predecessor paper calls the unity of the subjective and the objective: relational quality genuinely rises, and both partners sincerely experience it as rising, so that the improvement is not an artifact of adaptive preference \parencite{sen1999development,nussbaum2000women}. No one is deceived at the point of delivery. The question is whether, so stipulated, anything remains wrong; and if something does, what kind of thing it is, in whose jurisprudence it is legible, and what could cure it.

\subsection{Method and route}
\label{p02:sec:intro-method}

The paper proceeds dialectically, and announces its route in advance because the route is itself a claim. We begin where the worry begins, in the Hegelian-Marxian tradition, because that tradition possesses the only developed conceptual instrument for the intuition at stake: that a process can be defective while its product is good. Section \ref{p02:sec:hegel} establishes the frame: from Hegel, mediation as such is innocent, externalization (Ent\"au{\ss}erung) is the normal movement of spirit and of love, and the candidate for wrongness must therefore be located not in mediation but in a specific failure of its movement, externalization without return. Section \ref{p02:sec:marx} develops the Marxian antithesis in full: the fourfold structure of alienation in the 1844 manuscripts transposed to relational labor, with three further instruments (the Grundrisse's fragment on machines, commodity fetishism, and the distinction between formal and real subsumption) imported as structural transplants under an explicit methodological declaration. Section \ref{p02:sec:thesis} then gives the opposing thesis its strongest form, in the universalist grammar of modern law: instruments do not break causation, mediated sincerity is everywhere recognized, no harm grounds no action, and a reductio argument (the search engine, the friend, the relationship book itself) destroys every categorical version of the claim that mediation denatures intimacy. Section \ref{p02:sec:buddhism} deepens the thesis beyond what law can say, through the Buddhist analysis of dependent origination, and finds inside Buddhism its own strictly non-delegable core. Section \ref{p02:sec:synthesis} constructs the synthesis as a piece of doctrine: a default rule of instrumental imputation bounded by a criterion-governed exception domain modeled on the legal category of strictly personal acts, with a two-layer architecture of grounds and operationalization, a dynamic fifth criterion on which three unrelated traditions are shown to converge, and a touchstone case. Section \ref{p02:sec:poleconomy} then submits the synthesis to the one examination the dialectic has not yet conducted in the examiner's own voice, that of political economy: intimate acts are analyzed as layered phenomena bearing an allocative and a constitutive face; the jurisdiction of every allocative grammar is shown to end at the boundary between the faces, not at the boundary of the intimate domain; the specific danger of AI mediation receives its political-economic name, distributive reduction; and the mechanism of attribution-baseline drift is derived inside allocative grammar itself, as a deflation of the gesture's value-form. Three further tribunals then complete the tour, each summoned by the residue of the last. Section \ref{p02:sec:exploitation} conducts an eliminative experiment on exploitation theory: capital is removed from the scenario by construction, the self-built system examined in single-builder and joint-builder configurations, and a wrong is isolated that is insensitive to ownership configuration, surviving even the freely associated producers at the scale of two; the section then converts the eliminative result into a constructive program, generative justice practiced at the scale of the dyad, in which the paper's scattered remedies are unified as one four-layer circulation of value back to its generators. Section \ref{p02:sec:power} develops the single-builder residue, domination without exploitation, into a theory of the \emph{architectural constitutionalization of intimacy}: the mediating layer as a constitution drafted by one party, operating as third-dimensional power in the form of the good, against which subsequent consent is habitation rather than authorship, and whose only symmetric form is constitution practiced as a verb. Section \ref{p02:sec:feminism} convenes the tribunal longest owed the floor: it settles the paper's debts to feminist theory, shows the cognitive budget and the constituent power to be gendered in fact, distinguishes the redistribution of attending from its counterfeit, automated substitution, and supplies, from the ethics of care, the affirmative theory of the constitutive face that every prior section had used negatively. Section \ref{p02:sec:implications} returns the answer to the predecessor paper and draws the consequences for the boundary of consent in intimate life. Section \ref{p02:sec:limits} states limitations.

Two stipulations govern scope. First, the analysis concerns the structure of the mediating act, not the political economy of any provider. The system under examination may be assumed, throughout, to be self-built, self-hosted, and free of any third party's commercial interest; the paper's claims are designed to survive, or fail, under that assumption, because a critique that dissolves when the ownership form changes would not touch the question actually asked.\footnote{This stipulation does not deny that commercial deployment can aggravate matters; it brackets aggravation in order to isolate the base structure. See \S\ref{p02:sec:marx-method} for the corresponding methodological declaration on the Marxian side.} Second, the paper is jurisprudential and philosophical throughout: it derives no product guidance, no configuration advice, and no design recommendations, and it treats the predecessor paper's application only as a source of cases.

\subsection{Claimed contributions}
\label{p02:sec:intro-contributions}

The paper claims twelve contributions, stated here so that the reader may hold the argument to them. (1) The \term{contextual integrity of attribution}: a mirror-image transposition of \textcite{nissenbaum2009privacy}, on which what context-relative norms govern is not the flow of information but the flow of authorship credit for cognitive labor. (2) A theorization of the \term{supply chain of thoughtfulness}, naming with Marxian precision (\S\ref{p02:sec:marx-fetish}) the mechanism by which authorship dilutes and misreads along a multi-link chain of cognitive labor. (3) A doctrinal transplant: the extension of the category of strictly personal acts to intimate cognitive-caring acts, with the consequence that imputation fictions fail within the extended category. (4) The \term{process-constitutes-outcome} argument: a non-deontological explanation of why consequentialist evaluation does not so much lose as fail to find its object in this domain. (5) The identification of an \term{ironic form of alienation}: estrangement of person from person occurring under the description, and by means, of improving their relationship. (6) A two-wrongs architecture separating misattribution (curable by transparency) from process alienation (not curable by consent), which yields a principled answer to whether intimate life contains an unwaivable core. (7) The analysis of intimate acts as \term{layered phenomena} and the consequent thesis that the jurisdiction of allocative grammars ends at the boundary between an act's faces, not at the boundary of the intimate domain, with \term{distributive reduction} naming the characteristic transgression: the optimization of an act's allocative face that silently substitutes its constitutive face. (8) A political economy of the cognitive budget, comprising the budget/expenditure distinction, the identification of a managerial migration of attending (a \term{Taylorism of love}), and the \term{deflation of the gesture's value-form}, which supplies the heretofore missing mechanism of attribution-baseline drift. (9) An eliminative test across ownership configurations (commercial platform, single builder, joint builders) establishing the existence of an \term{ownership-insensitive wrong}, and thereby the exact jurisdiction of exploitation theory in the intimate domain. (10) The \term{architectural constitutionalization of intimacy}: a transposition of the third dimension of power and of architecture-as-regulation into the dyad, on which a self-built mediating layer is a unilaterally drafted constitution against which subsequent consent is habitation rather than authorship. (11) The distinction between the \term{redistribution} of attending and its automated \term{substitution}: the gendered form of distributive reduction, on which a standing equality claim can be answered in counterfeit currency. (12) The unification of the paper's remedies, self-hosting, attribution hygiene, scaffolded return, and co-constitution, as the four-layer practice of generative justice in the intimate dyad, including the identification of the recipient as the uncredited co-generator of the gesture directed at her, and the derivation of the dynamic criterion, a fourth time, from the recursion principle.

% =====================================================================
\section{The Dialectical Frame: Mediation, Externalization, and Return}
\label{p02:sec:hegel}
% =====================================================================

Any tribunal convened against mediation must first survive a Hegelian objection so fundamental that, unanswered, it ends the proceedings: there is nothing else but mediation. The opening movement of the \emph{Phenomenology} is the demonstration that immediacy is a myth; sense-certainty, which believes itself to grasp the pure singular \emph{this}, discovers that even its barest deixis is already universal, already linguistic, already mediated \parencite{hegel1807phenomenology}. What holds for knowing holds, in Hegel's early writings and throughout the mature system, for love: love is precisely the experience of finding oneself only in another, a unity that exists in and through otherness, never prior to it. A lover who has read poems, inherited a mother tongue's vocabulary of tenderness, absorbed an upbringing's grammar of attention, is not thereby a diminished lover; he is the only kind there is. If the indictment of the tender tip is that a third term has entered between lover and beloved, Hegel answers that the third term was always already there, and that its name is spirit.

\subsection{Ent\"au{\ss}erung and Entfremdung}
\label{p02:sec:hegel-distinction}

The frame, however, contains the resources for its own qualification, and the qualification is the engine of this paper. Hegel distinguishes, in effect if not always in terminology, between externalization and estrangement. \term{Ent\"au{\ss}erung}, externalization, is the normal and necessary movement of spirit: spirit becomes actual only by going out of itself, embodying itself in works, institutions, words, and gestures, and then \emph{returning to itself} out of this otherness, enriched, recognizing the externalized object as its own. The movement has three beats, not two: exit, dwelling-in-otherness, return. \term{Entfremdung}, estrangement, names the pathology specific to the third beat: externalization that does not come back, an objectification that stands over against its author as something alien, in which the author can no longer recognize, or no longer exercises, his own activity.

This distinction does two things for the present inquiry. Negatively, it disqualifies the most tempting form of the indictment. The mere fact that a gesture of love passed through a mediating instrument, a poem, a friend's phrase, a machine's suggestion, cannot constitute the wrong, since passage through otherness is the form of every actual love. Positively, it relocates the question onto a dynamic axis. The question is not \emph{whether} the lover's attending was externalized into an instrument, but whether the movement completes: whether what was externalized is reappropriated, taken back up into the lover's own growing capacity to attend, or whether it remains outside him, a standing alien power on which he comes to depend and against which his own capacity atrophies. We will state this formally in \S\ref{p02:sec:synthesis-dynamic} as the dynamic criterion; for now it suffices to fix the image. A child's training wheels are externalization with return: the capacity migrates into the child, and the apparatus falls away. A prosthesis that replaces a function, such that the function never migrates and the organ wastes, is externalization without return. The Hegelian frame predicts, what the static vocabulary of ``using a tool'' cannot even express, that one and the same act-type, accepting a machine's suggestion to express appreciation, may be innocent or estranged depending on the direction of motion of the capacity over time.

\subsection{What the frame leaves open}
\label{p02:sec:hegel-open}

Hegel supplies the grammar of the problem but not its verdict, for two reasons. First, the return movement in Hegel is ultimately guaranteed: the system's confidence is that spirit always comes home. The phenomena before us offer no such guarantee; whether a particular lover's outsourced noticing returns to him as cultivated attentiveness or hardens into dependence is an empirical and practical question, to be adjudicated case by case, which is to say, the Hegelian frame, honestly applied to finite life, issues not in a theodicy but in a standing tribunal. Second, Hegel's vocabulary does not yet distinguish the positions \emph{within} the relation: it speaks of spirit's self-estrangement, not of what the mediation does to the second person, the Roxane position, whose attribution of authorship is the hinge on which the predecessor paper's concept of proxy epistemic injustice turns. For the first deficiency we will need Marx, who took exactly this step: he de-theologized Ent\"au{\ss}erung, denied the guaranteed return, and built a structural account of the conditions under which externalization congeals into estrangement. For the second we will need, much later, the law of attribution (\S\ref{p02:sec:synthesis-criteria}). The order of the paper follows the order of these debts.

% =====================================================================
\section{The Marxian Antithesis: Alienation of the Relational Process}
\label{p02:sec:marx}
% =====================================================================

\subsection{Methodological declaration: a structural transplant}
\label{p02:sec:marx-method}

This section transplants the Marxian theory of alienation from waged labor to intimate cognitive-caring labor, and the transplant must be declared, because its legitimacy is contestable and the paper's no-overclaim norm requires the contest to be visible. In the 1844 manuscripts, alienation is analyzed within a determinate social form: the worker's process and product belong to another, labor is performed under compulsion, and the fourfold estrangement unfolds from that premise \parencite{marx1844economic}. The tender tip scenario lacks employer, compulsion, and extraction; under this paper's scope stipulation (\S\ref{p02:sec:intro-method}) it lacks even a commercial provider. What is transplanted is therefore not the political-economic context but the \emph{structure}: the separation of an agent from the process of an activity that is essentially his; the independence of this defect from the quality of the product; the resistance of this defect to consent. The transplant travels a well-established scholarly lineage, the philosophical-anthropological reading of the 1844 manuscripts developed by \textcite{fromm1961marx} and \textcite{ollman1971alienation} and renewed, in deliberately post-metaphysical form, by \textcite{jaeggi2014alienation}, on whose reformulation, alienation as a ``relation of relationlessness,'' a defective relation of the agent to her own activity and world rather than the loss of a substantial essence, this paper will principally rely. The reliance is doubly motivated: Jaeggi's version needs no doctrine of a fixed human essence, which immunizes the antithesis against an objection from the Buddhist quarter that would otherwise dismantle it (\S\ref{p02:sec:buddhism-anatta}), and it is congruent with the relational ontology of the predecessor framework, in which relations precede relata \parencite{huang2025relational}. Three further Marxian instruments, the fragment on machines, commodity fetishism, and the distinction between formal and real subsumption, are likewise imported as structures, with their capital-theoretic contexts bracketed; the bracketing is noted once here and not repeated.

One further declaration completes the methodological accounting, and it answers in advance an objection the paper will later raise against itself (\S\ref{p02:sec:pe-budget}). The Marxian corpus straddles, internally, the very grammatical line this paper will draw between allocative and constitutive analysis (\S\ref{p02:sec:thesis-jurisdiction}, \S\ref{p02:sec:poleconomy}). The mature analytic apparatus of the critique of political economy, value, surplus, exploitation as the misappropriation of a measurable surplus, speaks the allocative grammar, and the paper will use it only where that grammar holds jurisdiction (\S\ref{p02:sec:pe-deflation}). The 1844 theory of alienation speaks a different grammar altogether: its indictment is precisely that a constitutive activity, free conscious activity whose value is internal to its exercise, has been degraded into a mere means for producing external products; it is not an allocative critique but the theory of what happens when allocative grammar overreaches. The relation between these two Marxes is a contested field, and the paper declares its side rather than feigning a unanimous author: against Althusser's thesis of an epistemological break on which the humanist alienation theory is superseded \parencite{althusser1965for}, the paper stands with the lineage running from \textcite{lukacs1923history} through the Frankfurt School to \textcite{honneth2008reification}, on which the problem of alienation and reification is continuous through the corpus; and it notes that Marx's own polemic against the ``vulgar socialism'' that revolves around distribution \parencite{marx1875critique} licenses the reading on which his deepest critique was never a contribution to allocative grammar but a delimitation of it.

\subsection{The fourfold structure transposed}
\label{p02:sec:marx-fourfold}

The 1844 manuscripts distinguish four estrangements: of the worker from the product of labor; from the process of labor; from species-being; and of man from man \parencite{marx1844economic}. Their transposition to the tender tip is not uniform, and the non-uniformity is informative.

\subsubsection{From the product: the inapplicable first moment}
\label{p02:sec:marx-product}

The first estrangement transposes worst, and its failure to transpose is the first analytic gain. The gesture, the product, is not lost to the lover; it is delivered, received, and credited to him; if anything the product is hyper-appropriated, credited to him more fully than his contribution warrants. This inversion signals at the outset that the wrong, if any, will not be found at the product pole, which is exactly where consequentialist evaluation looks. The scenario stipulates a good product. Marx's first service to the inquiry is the observation, central to everything that follows, that \term{alienation is not a defect of the product}: alienated labor can and does produce excellent commodities. An analysis that inspects only the gesture and its effects inspects the one location where, by construction, nothing is wrong.

\subsubsection{From the process: the central moment}
\label{p02:sec:marx-process}

The second estrangement transposes exactly, and is the heart of the antithesis. For Marx, estrangement in the process means that the activity, while it occurs, is not the worker's own: he is at home outside it and outside himself within it; the living core of the activity has been transferred elsewhere, and what remains to him is execution. Transpose: the cognitive-caring labor of intimacy, noticing the lowered mood, remembering the offhand remark about the flowers, inferring that this week she needs words of appreciation rather than acts of service, is performed by the machine; what remains to the lover is the execution of a gesture whose cognitive interior he did not inhabit. The predecessor paper's formula, that the preciousness of ``he remembered the flowers I love'' lies in the fact that \emph{he} remembered, is the second estrangement stated from the recipient's side. Stated from the agent's side: the lover is present in his own courtship as the performer of another intelligence's findings. He is, in the idiom this paper will justify presently, at home outside his own attending.

\subsubsection{From species-being: a moment held under guard}
\label{p02:sec:marx-species}

The third estrangement, from species-being (Gattungswesen), free conscious activity as the human essential power, would seem to transpose grandly: if attending and loving belong to the activities in which the human being is most herself, their outsourcing is estrangement from the essence. The paper declines to lean on this moment, for two stated reasons. First, essence-talk is precisely what the Jaeggi reformulation was adopted to avoid, and what \S\ref{p02:sec:buddhism-anatta} will show cannot survive the doctrine of non-self. Second, and decisive on its own, the species-being argument in unguarded form proves too much: reading a book also ``outsources'' moments of free conscious activity, and the argument cannot distinguish the cases. It is therefore held under guard until the criteria of \S\ref{p02:sec:synthesis-criteria} are available to bound it; within those bounds it survives as rhetoric, not as an independent premise.

\subsubsection{From person to person: the ironic form}
\label{p02:sec:marx-person}

The fourth estrangement, of man from man, transposes with a torsion that this paper claims as a contribution. In Marx, the estrangement of each from his activity issues in the estrangement of each from the other; the social bond is hollowed from inside. In the tender tip, the hollowing occurs under the opposite description: the mediation is adopted \emph{in order to} bring the persons closer, and, by stipulation, it phenomenally succeeds. Yet structurally, a third intelligence now performs part of the knowing of each by the other; the relation's interior traffic, attention, interpretation, response, is partially re-routed through a node that is neither of them; person faces person across an inserted layer whose presence the receiving side does not, in the moment of reception, perceive. This is the \term{ironic form of alienation}: an estrangement of person from person enacted in the name, and with the felt phenomenology, of intimacy. Its irony is not decoration. It explains why the wrong is invisible to both first-person reports at once, and therefore why the unity of subjective and objective improvement, the strongest fact the thesis side possesses, cannot close the case: both reports issue from inside the very structure whose insertion is in question. The predecessor paper's I--It drift \parencite[\S 8.1.3]{huang2025toward}, sharpened: the drift here travels disguised as the I--Thou's own deepening.

\subsection{Three further instruments}
\label{p02:sec:marx-instruments}

\subsubsection{The fragment on machines: knowledge objectified, agent demoted}
\label{p02:sec:marx-fragment}

In the Grundrisse's fragment on machines, Marx describes the absorption of the worker's skill and knowledge into the machine system as objectified ``general intellect,'' with the worker re-positioned as the machinery's watchman and regulator \parencite{marx1857grundrisse}. Bracket fixed capital; keep the structure, and the description fits the present case with uncomfortable precision. The accumulated arts of intimate attention, how to read a mood, when a remembered detail will matter, the entire informal craft that lovers have transmitted as practice, is objectified into the suggestion engine; the lover is re-positioned as the executor of its findings, the watchman of his own courtship. The fragment contributes what the 1844 manuscripts alone do not: a specific account of \emph{deskilling}, the prediction that capacities which migrate into the apparatus atrophy in the agent. This is the materialist underwriting of the Hegelian return-failure (\S\ref{p02:sec:hegel-distinction}): externalization without return is not merely a spiritual misfortune but an observable trajectory of capacity loss, and it gives the dynamic criterion of \S\ref{p02:sec:synthesis-dynamic} its empirical content.

\subsubsection{The fetishism of the gesture}
\label{p02:sec:marx-fetish}

Commodity fetishism names the mechanism by which definite relations among persons assume the form of properties of things: value appears as an intrinsic attribute of the commodity, and the chain of labor behind it is occluded \parencite{marx1867capital}. Transpose the form, bracketing exchange: in the received gesture, a chain of cognitive labor, machine inference upon relational data, design of the suggestion, the lover's selection and execution, appears as a simple intrinsic property of the lover, his thoughtfulness. The predecessor paper's phrase, a supply chain of thoughtfulness whose origin the final recipient misreads, can now be given its theoretical name: the wrong at the reception pole is a \term{fetishism of the gesture}, the occlusion of a labor chain behind an apparent intrinsic quality. The naming is not ornamental. It imports Marx's central analytic lesson about fetishism: the misreading is not a private cognitive error of the receiver, to be cured by her being cleverer, but a \emph{structural} effect of the arrangement, which presents the gesture in a form that systematically invites the misattribution. This will matter doctrinally in \S\ref{p02:sec:synthesis-criteria}, where reasonable reliance, a structural and not a psychological standard, becomes the legal register of the same point.

\subsubsection{Formal and real subsumption: the book and the engine}
\label{p02:sec:marx-subsumption}

Marx distinguishes the formal subsumption of labor, in which an existing labor process is taken over but internally unchanged, from real subsumption, in which the process itself is reorganized around the apparatus \parencite{marx1864results}. Bracket subsumption \emph{under capital}; the depth-distinction remains, and it performs, from inside the Marxian register, work that will prove decisive against the reductio of \S\ref{p02:sec:thesis-reductio}. When the lover reads a book on love languages, external knowledge enters, but the labor process of attending remains intact and his: he reads, digests, judges when and whether and how to apply, and the noticing that triggers application is his own. This is formal subsumption at most: a new input to an unreorganized process. When an ambient engine monitors the relational data and issues the prompt, the process itself is reorganized: the \emph{noticing} link is now performed by the apparatus, and the lover's process has been restructured around reception of its outputs. This is the structural analogue of real subsumption. The book and the engine, which the reductio will treat as isomorphic links in identical supply chains, are thus distinguishable on Marxian grounds alone, by the depth at which the apparatus penetrates the process: input into the process, versus replacement of a constitutive link of the process. The reader should hold this distinction; it returns as the third criterion (\S\ref{p02:sec:synthesis-criteria}) in doctrinal dress.

\subsection{Consent does not waive structural alienation}
\label{p02:sec:marx-consent}

The antithesis must finally confront the strongest liberal reply: suppose both partners know, both agree, both prefer the mediated arrangement, and both judge the relationship better for it. Volenti non fit iniuria; where is the wrong? The Marxian answer is structural: the worker's consent to the wage bargain, however genuine, does not dissolve the alienation of the labor process, because alienation is a property of the arrangement, not of the parties' attitudes toward it \parencite{marx1844economic,jaeggi2014alienation}. Consent answers complaints of the form ``this was done to me against my will''; it does not touch defects of the form ``this arrangement separates the agents from their own activity,'' which persist with the arrangement whatever is felt about it. Modern legal orders themselves recognize the pattern: labor law and consumer law maintain unwaivable protections precisely because some structural positions ought not to be alienable even by their occupants' agreement. Whether intimate life contains such an unwaivable core, and what exactly falls inside it, cannot be settled from within the Marxian register, which supplies the concept but not the boundary; the question is remanded to \S\ref{p02:sec:synthesis-twowrongs} and \S\ref{p02:sec:implications-consent}, where the two-wrongs architecture will give it a determinate and, the paper will argue, principled answer: transparency cures what consent can reach, and what consent cannot reach is bounded, not unlimited.

\subsection{The antithesis, stated}
\label{p02:sec:marx-stated}

The Marxian case is now complete in outline. The wrong of the tender tip, if it is one, is not in the product (which is good), not in deception at delivery (there is none), and not in mediation as such (Hegel forbids that move). It is a process defect of a specific, nameable shape: the constitutive interior of an act of intimate attending is performed by an apparatus while the act continues to present itself, and to be received, as the agent's own; the agent is demoted to executor (fragment on machines), the receiver's misattribution is structurally invited rather than accidentally committed (fetishism of the gesture), the penetration reaches the process's constitutive links rather than its inputs (real subsumption), and the whole occurs under the description of relational improvement (the ironic form), so that the phenomenology of both parties testifies for the arrangement and cannot be cited against the indictment without circularity. And the defect, being structural, is not dissolved by consent. What the antithesis cannot yet say is why this indictment does not also convict the book, the friend, and the search engine; it has gestured (subsumption) but not yet answered. The thesis side will now press exactly there, and with the full authority of the universalist legal tradition behind it.
% =====================================================================
\section{The Universalist Thesis: Tools Do Not Break Causation}
\label{p02:sec:thesis}
% =====================================================================

The antithesis now meets the modern legal tradition at its most confident, and the encounter must be staged fairly: the thesis is to be given its strongest form, because a synthesis bought by weakening one's opponent is worthless. The universalist grammar of law, common to the civil and common law families in everything that matters here, offers three replies and one devastating argument.

\subsection{The instrumentality doctrine}
\label{p02:sec:thesis-instrument}

The first reply is doctrinal bedrock. Interruption of causation, novus actus interveniens, superseding cause, the negation of objektive Zurechnung, is reserved for intervening events of a particular character: free, informed, voluntary interventions of third parties, or extraordinary unforeseeable occurrences, which displace the original actor's agency as the explanation of the outcome \parencite{hart1985causation,roxin2006strafrecht}. An instrument that the actor himself selects, foreseeing and intending its contribution, is the opposite of such an event: it is the very medium of his agency, and imputation passes through it without resistance. The marksman's rifle does not interrupt the causation of the death; the accountant's calculator does not make the software the author of the audit. The intervening-cause register, applied honestly, classifies the suggestion engine as transparent: the lover chose it, foresaw its function, adopted its output by a fresh act of his own will, and executed the gesture himself. On this analysis there is not even a colorable interruption; the love expressed, however machine-assisted, is imputed to the lover, fully and without remainder.

It is worth pausing on what this does to the antithesis's most legal-sounding formulation. One might have argued (the predecessor session did argue, before discarding it) that the recipient's attribution is, in causal-legal terms, an objective misattribution: she assigns the outcome to the origin of a chain that was in fact interrupted. The instrumentality doctrine destroys this formulation on its own terrain. The chain was not interrupted; her attribution of the gesture to her partner is, as far as causation runs, \emph{correct}. If a wrong survives, it cannot be a wrong of causal misattribution. The paper accepts this result without reservation; the synthesis will have to, and will, locate the misattribution elsewhere than in causation (\S\ref{p02:sec:synthesis-criteria}, fourth criterion).

\subsection{Mediated sincerity is everywhere}
\label{p02:sec:thesis-mediated}

The second reply is institutional. Law does not merely tolerate mediated authenticity; it is built of it. The will drafted by counsel is the testator's last will; the oath rendered through an interpreter is the deponent's oath; the contract concluded through an agent binds the principal as if he had signed: qui facit per alium facit per se. The entire edifice of agency exists to ensure that interposition does not disturb imputation, because a legal order in which only unmediated acts counted as one's own would be unworkable and, more deeply, false to the mediated texture of all social action, the same point Hegel made metaphysically (\S\ref{p02:sec:hegel}) now made institutionally. The universalist's first instinct, confronted with an intermediary, is not interruption but \term{penetrating fiction}: treat the mediated act as the principal's own. Against this background, the claim that a machine's participation in composing a gesture denatures the gesture is not a refinement of legal thought but a revolt against its load-bearing walls.

\subsection{No harm, no action}
\label{p02:sec:thesis-noharm}

The third reply is jurisdictional. Doctrines of causal interruption exist to \emph{exculpate}: to sever a defendant from a harm. Here there is no harm; the stipulated facts contain an improved relationship and two sincerely flourishing persons. Even misrepresentation, the nearest nominate wrong, requires a false representation, reasonable reliance, and detriment; on the stipulated facts the second element is contestable and the third is absent. Law has no action for the impurity of a causal chain as such, and a jurisprudence that invented one would have invented a wrong without a victim. The thesis presses the point past doctrine into method: perhaps the felt residue that the antithesis elaborates is not a wrong at all but a sentiment, a romantic preference for artisanal cognition, deserving of respect as taste but not of enforcement as norm.

\subsection{The reductio: the search engine, the friend, and Chapman's book}
\label{p02:sec:thesis-reductio}

The fourth argument is the strongest, and the paper adopts it as its own weapon rather than treating it as an embarrassment. If mediation of intimate cognition denatures the act, the principle proves vastly too much. The lover who searches for the city's best restaurant has outsourced cognitive labor to an index built by strangers; the lover who consults a travel guide has outsourced destination-choice to its author; the lover who asks a mutual friend what gift she would treasure has outsourced a tract of partner-modeling to the friend. And, the decisive instance: the lover who reads the five-love-languages book \emph{itself} has acquired, from a static external artifact, precisely the heuristic that the suggestion engine applies; the book is, in the relevant sense, a tip generator with a page interface. If the engine's mediation denatures, the book's mediation denatures; no one believes the book's mediation denatures; therefore the categorical principle, mediation of intimate cognition denatures the act, is false. The reductio is not a debater's trick. It encodes the universalist tradition's deepest insight, that human action is irreducibly prosthetic, and it sets the formal constraint that any surviving version of the antithesis must satisfy: the surviving version must be a \emph{discriminating} principle, one that sorts the engine from the book by stated criteria, or it is nothing.

\subsection{What the reductio destroys, and what it leaves standing}
\label{p02:sec:thesis-aftermath}

Take stock with precision, because the entire synthesis hangs on the exact shape of the wreckage. Destroyed: every categorical claim of the form ``AI mediation of intimacy is as such denaturing, alienating, or wrong.'' Destroyed with it: the unguarded species-being argument (\S\ref{p02:sec:marx-species}), and any reading of the antithesis on which the wrong follows from mediation alone. Left standing, because the reductio never touched them: (i) the Marxian observation that process defects are independent of product quality, which the no-harm reply (\S\ref{p02:sec:thesis-noharm}) does not refute but simply declines jurisdiction over; (ii) the subsumption distinction (\S\ref{p02:sec:marx-subsumption}), which already discriminates the book from the engine and which the reductio's isomorphism claim quietly assumes away; (iii) the fetishism analysis of structurally invited misattribution, which concerns the form in which the gesture presents itself, a matter on which the causal correctness of the attribution (\S\ref{p02:sec:thesis-instrument}) is silent; and (iv) the legal tradition's \emph{own} counter-category, so far unmentioned by the thesis side because it is the thesis side's internal exception: the class of acts that law itself refuses to let pass through intermediaries. The question has therefore transformed, exactly as dialectic requires. It is no longer whether mediation is permissible, to which the answer is yes by default, but which features place a given mediation inside a wrongful exception domain, and whether the tender tip possesses them. That is a question the universalist tradition is not only able to hear but has, in its own house, already begun to answer.

\subsection{The jurisdiction of imputation: a diagnosis, and its necessary limit}
\label{p02:sec:thesis-jurisdiction}

Before leaving the legal register, the paper owes a diagnosis of what, exactly, the thesis side's victories amounted to, for the victories were real and yet the unease survives them, and an unease that survives correct doctrine is usually evidence of a question the doctrine was not built to hear. The imputation apparatus, causation, novus actus, agency, is a retrospective, event-individuating, assignment-issuing machine. Its built premises are three: that there is a severable event (an act, a harm, a result) to serve as the object of imputation; that the event is complete, the judge standing after it and looking back; and that the output is an assignment, this act is yours, or it is not. Within those premises the apparatus is sovereign, and \S\S\ref{p02:sec:thesis-instrument}--\ref{p02:sec:thesis-reductio} were its sovereignty exercised. But the question the predecessor paper kept open is not of that shape. Whether a relationship is being well constituted, whether its manner of being sustained is building or hollowing the persons who sustain it, is processual where imputation is event-shaped, prospective where imputation is retrospective, and constitutive where imputation is assignative: the acts at issue do not produce results external to themselves to which authorship could then be assigned; they \emph{are}, ongoingly, the relationship.\footnote{The classical anchor is Aristotle's distinction between poiesis, making, whose end is a product separable from the activity, and praxis, doing, whose end is internal to its exercise \parencite[VI.4--5]{aristotle2009nicomachean}. Imputation grammar is built for poiesis-shaped acts; attending is praxis; and the Marxian diagnosis of \S\ref{p02:sec:marx} can be restated in exactly these terms, as the theory of what occurs when praxis is forcibly processed as poiesis.} The instrumentality doctrine, asked about the tender tip, answered correctly, and answered a question nobody had asked: imputation was never in doubt; formation was.

The diagnosis, however, must immediately be bounded, or it destroys more than the paper can afford, including the paper's own title. It would be false, and self-refuting, to conclude that legal grammar is unfit for the intimate domain as such. The line between the two question-types does not run between law and intimacy; it runs \emph{through intimate life itself}, and indeed, as \S\ref{p02:sec:poleconomy} will show, through every individual intimate act. Intimate life is dense with genuinely event-shaped, assignable matters, deception practiced, trust breached, harm inflicted, over which retrospective imputation holds full and beneficial jurisdiction: the law of divorce, of intimate wrongs, of misrepresentation operates there without category error, and will be needed again the moment the layered analysis of \S\ref{p02:sec:pe-layered} identifies genuinely allocative phenomena inside the home. Nor is the boundary an external imposition upon law by philosophers: law knows it. The relational contract theory of \textcite{macneil1980new} is an internal legal demonstration that the discrete-transaction paradigm cannot process ongoing relations and that a different normative grammar is required for them; the restraint traditions of family law, the reluctance of courts to adjudicate the interior conduct of intact marriages, encode the same self-knowledge institutionally; and the category of strictly personal acts, on which the synthesis will be built, is nothing other than the constitutive grammar's enclave \emph{inside} the law of acts. The thesis side's victory is therefore to be entered in the record precisely: complete within the jurisdiction of imputation, silent beyond it, and the silence is not a defect of law but the shape of its competence. What remains for the synthesis is to say where the jurisdiction ends, and \S\ref{p02:sec:thesis-aftermath}'s transformed question can now be stated in its final form: not whether mediation is permissible, nor even which mediations fall in an exception domain, but \emph{which face of a layered act a given mediation operates upon}.

% =====================================================================
\section{The Buddhist Examination: Dependent Origination and the Act That Was Never Unmediated}
\label{p02:sec:buddhism}
% =====================================================================

Before the synthesis, the thesis receives reinforcement from a tradition that owes nothing to Roman law or to Hegel, and the reinforcement arrives carrying, unasked, its own exception domain. The convergence between what this section finds and what \S\ref{p02:sec:thesis} found is itself a datum, and the paper will treat it as one (\S\ref{p02:sec:buddhism-convergence}).

\subsection{No first link: the tip as one condition among conditions}
\label{p02:sec:buddhism-conditions}

On the analysis of dependent origination (pa\d{t}icca-samupp\=ada), no act arises from an unconditioned interior. Every act of attending that a lover has ever performed was already constituted by conditions without limit: the books read, the language inherited, the sensibility formed by an upbringing, the friend's remark decades ago. The ``unmediated act of pure personal attending,'' which the alienation indictment treats as the displaced original, is on this analysis not displaced but nonexistent; there never was such an act to be displaced. The machine's suggestion is one further condition entering the conditioned stream, of no different ontological standing than the remembered poem. And moral weight has a precise location in this picture: \emph{cetan\=a}, volition, is karma \parencite[AN 6.63]{bodhi2012numerical}. The volitional sequence, the intention that turns toward the machine and asks; the discernment that receives and evaluates its answer as good or bad; the decision to act or refrain; the quality of heart in the acting, occurs entirely in the lover and is outsourced at no point. Where the universalist lawyer says the tool does not interrupt the causal chain, the Buddhist analysis says something stronger: there is no proprietary causal chain there to interrupt, only conditions and, threading them, volitional acts that remain ineluctably the agent's own. The doctrine of novus actus cannot even be formulated, for it presupposes a self-standing causal authorship that conditions never had.

\subsection{The deeper cut: non-self against the metaphysics of alienation}
\label{p02:sec:buddhism-anatta}

The same analysis now turns on the antithesis, and cuts deeper than law could. An alienation indictment of classical form requires something to be alienated \emph{from}: an essence, a species-being, an authentic self whose proper activity has been estranged. The doctrine of non-self (anatt\=a) denies that any such substantial owner exists; there is no svabh\=ava-bearing subject behind the stream of conditioned acts whose essential property the attending could be. Pressed without preparation, this dissolves the third estrangement entirely and threatens the rest. The paper, however, prepared (\S\ref{p02:sec:marx-method}): the Jaeggi reformulation on which the antithesis was built defines alienation not as separation from an essence but as a defective \emph{relation}, a relation of relationlessness between an agent and her activity, and a relational defect requires no substantial self, only the relation. Indeed the three vocabularies now visibly align: dependent origination, the predecessor framework's relational ontology, and Jaeggi's de-essentialized alienation are all assertions that the relation is prior to the relata. The antithesis survives the anatt\=a objection, but only in the post-metaphysical form this paper gave it from the start, and the survival is the second independent confirmation (after \S\ref{p02:sec:marx-method}) that the form was forced, not chosen.

\subsection{Buddhism's own non-delegable core}
\label{p02:sec:buddhism-core}

\subsubsection{Practice cannot be performed by another}
\label{p02:sec:buddhism-practice}

Yet the tradition that so thoroughly acquits mediation maintains, at its center, a category of the strictly personal more austere than anything in the civil codes. Purity and impurity are one's own; no one purifies another \parencite[v.~165]{buddharakkhita1985dhammapada}; be islands unto yourselves, refuges unto yourselves \parencite[DN 16]{walshe1995long}. The path admits kaly\=a\d{n}amittat\=a, good friendship, teachers, texts, and expedient means in abundance, but the walking of it is non-transferable without residue. And the faculty whose cultivation constitutes the path's core discipline is precisely \emph{sati}, attention, present-moment noticing. Herein the Buddhist register reframes the tender tip more exactly than either prior register could: what the ambient engine performs on the lover's behalf is not a task but a \emph{practice-moment}, an occasion for the cultivation of attentiveness, and practice-moments are the one currency the tradition holds to be strictly personal. Delegating the noticing is not theft from the beloved (the gesture arrives; she flourishes); it is, in the first instance, a forfeiture in the agent's own training, the atrophy that the fragment on machines predicted (\S\ref{p02:sec:marx-fragment}) now redescribed as the stunting of a cultivable faculty. Tronto's placement of attentiveness as the first moment of care \parencite{tronto1993moral} receives here an independent, non-Western grounding: attention is first not because care theory ranks it so, but because it is the faculty in which the entire ethical path is trained.

\subsubsection{The raft and clinging: the dynamic criterion, second derivation}
\label{p02:sec:buddhism-raft}

The simile of the raft \parencite[MN 22]{nanamoli1995middle} supplies the tradition's test for expedient means: the dhamma itself is a raft, for crossing over, not for holding onto; the worth of an instrument lies in its trajectory toward its own dispensability. An aid that is used and outgrown, the suggestion that teaches the lover to notice until the engine falls silent from disuse, is a raft, and blameless. An aid that is clung to, up\=ad\=ana, such that reliance deepens and the native faculty recedes, has become a fetter under the description of a vehicle, and clinging is, in the twelve links, exactly the pivot at which conditioned contact rolls onward toward bondage. The reader will recognize the structure: this is the Hegelian distinction between externalization-with-return and externalization-without-return (\S\ref{p02:sec:hegel-distinction}), derived a second time, from premises sharing no genealogy with the first derivation. Training wheels and prosthesis; raft and fetter. Two traditions, asked independently, answer that the moral character of the mediation is not legible in the act-token but only in the direction of motion of the agent's capacity over time.

\subsubsection{Right speech and the invited misreading}
\label{p02:sec:buddhism-speech}

Finally, the tradition's speech ethics reaches the reception pole. The discipline of right speech extends beyond the avoidance of false assertion to the non-cultivation of false impressions in others; a course of conduct that foreseeably and systematically leads another to form an erroneous belief stands in the penumbra of false speech even if no individual utterance is false. The arrangement in which gestures arrive bearing the structurally invited misreading of their authorship (\S\ref{p02:sec:marx-fetish}) sits in exactly this penumbra. Buddhism thus independently generates a concern with \emph{attribution}, the fourth of the criteria to come, just as it independently generated the dynamic criterion and the non-delegable core.

\subsection{The convergence as a finding}
\label{p02:sec:buddhism-convergence}

The paper now possesses a result no single section could have produced. Three traditions without common ancestry in the relevant respects, the universalist legal grammar, the Hegelian-Marxian dialectic, and the Buddhist analysis of conditioned action, have each delivered the same tripartite shape: a default acquittal of mediation as such; a bounded domain of the strictly personal in which the acquittal fails; and a dynamic test (return or no return; raft or fetter) for the cases the boundary leaves uncertain. In the method of the predecessor paper, which assigns jurisdictions to traditions and trusts no single one with the whole \parencite{huang2025toward}, a convergence of this kind is the strongest species of evidence available: the synthesis about to be constructed is not the verdict of a school but the overlapping consensus of three. The construction may now proceed on doctrinal ground, with the other two traditions standing as independent sureties for its shape.
% =====================================================================
\section{The Synthesis: Strictly Personal Acts and the Failure of Imputation Fictions}
\label{p02:sec:synthesis}
% =====================================================================

\subsection{The doctrinal anchor: h\"ochstpers\"onliche Handlungen}
\label{p02:sec:synthesis-anchor}

The universalist tradition's imputation fictions, agency, instrumentality, penetration, rest on a tacit premise: that the normative value of the imputed act does not depend on who personally performs it. For conveyances, payments, and declarations of commercial will, the premise holds and the fictions do their indispensable work. But the tradition itself has always maintained a counter-category of acts for which the premise fails, the strictly personal acts, h\"ochstpers\"onliche Rechtsgesch\"afte, acts intuitu personae, and the membership of the category is instructive. Marriage must be contracted by the parties personally and in simultaneous presence (so, expressly, \S 1311 BGB); a will can be made only personally, never through a representative (\S 2064 BGB); Chinese law requires both parties to appear in person for marriage registration (Civil Code of the PRC, art.\ 1049); Japanese doctrine treats acts of family status (\textit{mibun k\=oi}) as resistant to agency in principle; the vote may not be cast by another in one's name; and criminal liability for status offenses does not transfer.\footnote{The anchor texts, verified against current sources: \S 1311 S.~1 BGB: ``Die Eheschlie{\ss}enden m{\"u}ssen die Erkl{\"a}rungen nach \S~1310 Abs.~1 pers{\"o}nlich und bei gleichzeitiger Anwesenheit abgeben''; representation and even declaration through a messenger are excluded, and recent German decisions have held video-link participation insufficient for the simultaneous presence the provision requires (so OVG Berlin-Brandenburg and, in 2024, the BGH on a ceremony conducted by a video-linked Utah official; >> VERIFY docket numbers for pin cites <<). \S 2064 BGB: ``Der Erblasser kann ein Testament nur pers{\"o}nlich errichten'' (formelle H{\"o}chstpers{\"o}nlichkeit; representation void, mere professional advice with final decision by the testator permitted, a statutory ancestor of this paper's pulled-means/pushed-noticing distinction); see also \S 2065 BGB, excluding determination of the disposition's content by a third party. Civil Code of the PRC, art.~1049: {\cjkfont 要求结婚的男女双方应当亲自到婚姻登记机关申请结婚登记} (both parties must apply in person), reiterated in art.~7 of the 2025 Marriage Registration Regulations ({\cjkfont 婚姻登记条例}, State Council Order No.~804). Japanese doctrine holds that acts of family status, marriage, divorce, acknowledgment, and wills, ``do not admit agency'' ({\cjkfont 身分行為は代理に親しまない}), the principal's true intention being treated as absolute (>> VERIFY treatise pin cite, e.g.\ a standard family-law commentary <<). Common law and French positions remain to be surveyed in the comparative section of a subsequent draft, reusing the author's six-jurisdiction methodology from prior work on AI-extensible products.} The category's rationale, articulated or implicit across jurisdictions, is uniform: there exists a class of acts whose normative force \emph{consists in} their personal performance, such that a substituted performance is not a defective instance of the act but a different act altogether. For these, the imputation fiction is not rebutted; it never engages.

The category, moreover, is candid about its own conventionality, and the candor must be faced rather than hidden, because it contains the strongest objection to what follows. Proxy marriage exists: historically among royalty, currently for deployed soldiers in a handful of jurisdictions. If even marriage's personalness is historically variable, is the category anything more than congealed social expectation? The paper's answer, developed below as the two-layer architecture (\S\ref{p02:sec:synthesis-architecture}), is that the objection proves exactly the right amount: the \emph{operationalized boundary} of the category is indeed conventional and moves with social expectation (which is why criterion four is formulated in terms of expectation baselines), while the \emph{grounds} of the category, the features in virtue of which personal performance is constitutive for some acts and not others, are not conventional, and it is the grounds that criteria one through three articulate. Proxy marriage is thus not a counterexample but a demonstration of the architecture: a jurisdiction that admits it has moved the conventional boundary, under pressure of war or distance, without anyone supposing that the grounds (the wedding of \emph{this} person to \emph{this} person) had changed; and the institution remains marked everywhere as exceptional, boundary-stretching, in need of special justification, which is how a conventional surface behaves when it is anchored to non-conventional grounds.

\subsection{The claim}
\label{p02:sec:synthesis-claim}

The synthesis can now be stated in one sentence and defended in the rest of the section. \emph{Intimate cognitive-caring acts belong, or ought to be recognized as belonging, to the category of strictly personal acts; within this category the imputation fictions fail, not because the instrument interrupts causation (it does not), but because acts of this class do not admit substitutable links; and the boundary of the class is given not by the bare presence of mediation, which the reductio has made untenable, but by four criteria, beneath which lie graded grounds and above which operates a dynamic test.} The claim is doctrinally conservative in form, default imputation plus an enumerated exception domain is precisely how law has always handled agency, and it is this conservatism that makes it strong: the paper asks the universalist tradition to extend a category it already maintains, not to adopt a foreign one.

\subsection{The four criteria}
\label{p02:sec:synthesis-criteria}

The criteria sort mediations, not technologies; they apply to the friend and the book as readily as to the engine, which is what the reductio demands.

\begin{enumerate}[leftmargin=2em, itemsep=0.35em]
\item \textbf{Object of mediation: world-knowledge or the other's mind.} The search engine mediates facts about the world, which restaurant is good, facts public in kind and owned by no relationship. The tender tip mediates an inference about a particular person's interior, that she has lately lacked words of appreciation, which is theory-of-mind labor, the cognitive core of intimate knowing. Consulting the world's knowledge of the world outsources nothing of the knowing-of-her; consulting an engine's model of her outsources a tract of exactly that.
\item \textbf{Provenance of the inference: public or relation-interior data.} The index does not know the beloved; the engine's inference runs on the sediment of the relationship itself, moods, diaries, remembered remarks, data generated inside the intimate context and governed, on \textcite{nissenbaum2009privacy}'s analysis, by that context's norms. A mediator operating on relation-interior sediment is not adjacent to the relationship but inside it, performing on the relationship's own materials the interpretive labor that belonged to its members.
\item \textbf{Direction of initiative: pulled means or pushed noticing.} This is the subsumption distinction (\S\ref{p02:sec:marx-subsumption}) in doctrinal dress, and the Tronto-sati point (\S\ref{p02:sec:buddhism-practice}) in legal register. When the lover, having already formed the intention to delight her, pulls an instrument toward chosen ends, intention, attention, and initiation remain his; the instrument fills in means, and means are the natural home of mediation. When an ambient system initiates, pushing ``notice her now'' upon an agent who was not noticing, the outsourced element is not means but the \emph{inception of attending itself}, the first moment of care, the practice-moment; and the inception is the link that the grounds mark as constitutive. (Herein the predecessor paper's design constraints, opt-in and infrequent, are revealed as having tacitly honored this criterion: they were devices for keeping initiative on the human side.)
\item \textbf{Attribution expectations: the contextual integrity of attribution.} No recipient supposes her partner personally surveyed the city's restaurants; the use of an index lies within the standing social baseline of expected mediation, and her attribution of the evening to his care is undisturbed by truth. But ``he noticed I was low; he remembered what I said in passing'' is attributed, under the present baseline, to personal attending, and a structure that produces such gestures from machine noticing invites, systematically and by its form (\S\ref{p02:sec:marx-fetish}), an attribution the facts do not honor. The norm violated is the mirror image of Nissenbaum's: not the contextual integrity of information flow but \term{the contextual integrity of attribution}, the context-relative norms governing where authorship credit for cognitive labor may flow. Its legal register is reasonable reliance: a structural baseline, not a psychology, which is why the fetishism analysis and the reliance doctrine, arriving from opposite traditions, state one standard.
\end{enumerate}

\subsection{The two-layer architecture: grounds and operationalization}
\label{p02:sec:synthesis-architecture}

The criteria are not four parallel switches, and treating them so would invite two fatal misreadings: that criterion four, being the most legally tractable, simply absorbs the others (whereupon the category collapses into convention and the proxy-marriage objection wins); or that criteria one through three operate as freestanding prohibitions (whereupon the friend's kind counsel and the couples therapist are absurdly convicted). The architecture is two-layer. Criteria one through three are \emph{grounds}: they articulate the features in virtue of which an act is a candidate for the strictly personal category at all, mediation reaching the other's mind rather than the world, working the relation's interior sediment, replacing inception rather than means. They are graded, not binary, and their joint satisfaction in high degree is what makes substituted performance a different act rather than an assisted one. Criterion four is the \emph{operationalization}: the legally administrable surface where the grounds meet a community's standing baseline of expected mediation, and the only layer at which formal doctrine should act, precisely because baselines are public, provable, and revisable, while grounds are matters of philosophical articulation. The layers interact in one direction over time: as the grounds' satisfaction becomes commonly known for a practice, the baseline shifts, and the operationalized boundary follows, which is the lesson of proxy marriage generalized, and the reason the doctrine here proposed will not be stranded by technological change.

The architecture also disposes of the cases that would otherwise embarrass the criteria. The mutual friend who volunteers ``her birthday is near; she mentioned wanting that book'' trips, on paper, criteria one through three; yet the standing baseline of every human community has always included the counsel of friends, matchmakers, and confidants within expected mediation, so criterion four is not violated and no misattribution is invited; while the residual process-question (did the lover let friends do all his noticing?) is exactly the kind that the grounds mark as graded and the dynamic criterion (next subsection) adjudicates. Cyrano, by contrast, is convicted, and should be: the baseline of the balcony does not include a hidden poet, the attribution invited is total, and the play knew it. The wrong this paper articulates is therefore not specific to artificial intelligence, the briefing's instinct to open with Rostand was sound, and the honest statement of AI's contribution is quantitative and structural: it renders Cyrano's arrangement scalable, ambient, perpetual, and available within every relationship at zero marginal cost, thereby pressing on the baseline itself, a pressure whose normative assessment is taken up at \S\ref{p02:sec:implications-consent}.

\subsection{The dynamic criterion: externalization-and-return}
\label{p02:sec:synthesis-dynamic}

Within the domain the grounds demarcate, a residue of cases remains in which the static features genuinely underdetermine judgment, the same suggestive practice that entrenches one lover's dependence educates another into attentiveness. For these the paper promotes to criterial status the test that two traditions independently generated: \emph{whether the externalization returns}. If the mediation operates as training wheels, raft, formal subsumption, if the capacity migrates into the agent, observable as a rising autonomy of noticing and a falling reliance on the apparatus, the movement is the ordinary, innocent movement of spirit through otherness, and no alienation has occurred whatever the static criteria scored. If it operates as prosthesis, fetter, real subsumption, reliance deepening, the native faculty receding, the externalization has not returned, and the arrangement stands convicted on the only timescale at which process defects are visible at all. The dynamic criterion explains, what no static analysis could, why a single act-token of accepting a suggestion is morally illegible in isolation, and it binds the synthesis back to the predecessor paper's dialectical-positivist method: the question ``does this measure and describe the relationship, or has it begun to replace it'' is asked anew at every site, and is never finally settled \parencite{huang2025toward}.

\subsection{Two wrongs, two remedies}
\label{p02:sec:synthesis-twowrongs}

The machinery now resolves the tension that \S\ref{p02:sec:marx-consent} left standing, between the Marxian axiom that consent does not waive structural alienation and the manifest curative power of disclosure upon misattribution. The tension dissolves because there are two wrongs, of different kinds, with different victims and different cures. The first is \term{misattribution}: an epistemic-relational wrong at the reception pole, the invited misreading of the supply chain of thoughtfulness, the violation of the contextual integrity of attribution; its primary victim is the recipient; and it is curable by transparency, where both partners know the practice, the attribution is corrected at the source, the fetish is, as it were, labeled, and this wrong does not arise. The second is \term{process alienation}: a structural wrong at the agency pole, the replacement of constitutive links of intimate attending, the relation of relationlessness between the lover and his own loving; its victims are the agent and the relation itself; and it is \emph{not} curable by consent, for the Marxian reason that it is a property of the arrangement and not of attitudes toward it, but it is \emph{bounded}, arising only within the domain the grounds demarcate and only on the failing branch of the dynamic criterion. Consent, in short, cures everything that was ever about deception and nothing that was ever about structure; and the structure's reach is not unlimited but criterial. This is the paper's answer to the question whether the liberal or the Marxian view of consent governs intimate life: each governs its own wrong.

\subsection{The touchstone: the self-built system}
\label{p02:sec:synthesis-touchstone}

The synthesis is best exhibited on the hardest clean case, a system built by the lover himself, hosted by himself, serving no interest but the couple's, the case this paper's scope stipulation has held in view throughout. The analysis does not acquit it wholesale, and does not convict it wholesale; it cuts through the middle of it, which is the demonstration that the blade is sharp. The module in which the lover has objectified \emph{his own past attending}, the anniversaries he entered, the flowers he once noticed and recorded, and which later returns these to him, is externalized memory completing its return: structurally a diary with an alarm, well inside the baseline of expected mediation, initiating nothing that the lover's past self did not initiate. It passes all four criteria and the dynamic test, and the analysis finds in it nothing to forgive. The module in which an inferential engine runs on the relation's live interior sediment and initiates findings about her present mind, however lovingly built, scores against grounds one through three, presses against the attribution baseline, and stands or falls thereafter on the dynamic criterion alone: raft if it trains him out of needing it, fetter if it trains him into needing it. Same author, same hardware, same data, same intention, opposite verdicts across two modules: the analysis is not about AI, not about ownership, not about commerce, and not about technology; it is about which links of a strictly personal act have been substituted, and in which direction the capacity is moving. That is the resolution the question deserved, and the only kind, this paper has argued, that it admits.

% =====================================================================
\section{The Political Economy of the Cognitive Budget}
\label{p02:sec:poleconomy}
% =====================================================================

The synthesis must now survive an examination it has so far postponed, and the examination is mandated by the paper's own honesty norms. The constitutive analysis of \S\ref{p02:sec:synthesis} would be worthless if it were purchased by denying the allocative realities of intimate life, and the denial would be doubly false: false to the facts, since attention is scarce in the strictest sense and the tender tip exists \emph{because} it is scarce (were attention inexhaustible, no machine would need to husband it); and false to the predecessor paper, whose household modules, ledgers, schedules, the deliberate offloading of factual memory, are allocative through and through, were defended as such, and work \parencite{huang2025toward}. Romantic purism, the doctrine that love lies outside economics, is not available to this paper and is not true. This section therefore grants political economy full jurisdiction over everything it can hear, and demonstrates a result more useful to the synthesis than any exemption: that allocative grammar, speaking strictly for itself, describes the entire outline of the danger, while the final account of why the danger is a wrong, and not an efficiency, remains where \S\ref{p02:sec:synthesis} placed it.

\subsection{Scarcity without transferability: budget and expenditure}
\label{p02:sec:pe-budget}

Begin with the definition that any account of intimate cognition must face. Economics, on Robbins's canonical formulation, studies human behavior as a relationship between ends and scarce means which have alternative uses \parencite{robbins1932essay}; and cognition is the paradigm of such a means. Attention is finite; in Simon's formula, a wealth of information creates a poverty of attention, and the poverty necessitates allocation \parencite{simon1971designing}. The hour of attending given to the beloved is an hour not given to the manuscript; the couple's evening is a budget line. Nothing in this paper denies, and much in the predecessor paper affirms, that intimate life is a site of real cognitive scarcity and therefore of real allocative problems, to which allocative reasoning is the correct response.

The precision the paper requires is a distinction \emph{within} the economics of attention, between two operations that a budget admits. The first is \term{intra-personal allocation}: the agent's distribution of his own finite cognitive budget across competing claims, work, household, relationship, rest. This is scarcity in Robbins's full sense, and the agent's economizing over it is ordinary, legitimate, and often obligatory rational conduct. The second is \term{inter-subjective substitution}: the replacement of the budget's expender, another intelligence spending in the agent's stead. The constitutive analysis of \S\ref{p02:sec:synthesis} concerns only the second operation, and the reason can now be stated as a boundary theorem on Robbins's definition rather than as a romantic exception to it. Constitutive attending does not fall outside the definition of the economic; it falls inside it as a special subclass: a scarce means \emph{whose use-value is inseparable from the identity of its expender}. The definition's ordinary apparatus, opportunity cost, efficient assignment, gains from trade, presupposes that means are transferable between expenders without loss of the end served; the presupposition holds for money, for hours of housework, for the booking of the restaurant, and fails, not contingently but constitutively, for the attending whose entire yield to its recipient is \emph{that he attended}. The intimate domain is not extraterritorial to economics; it contains a class of goods for which the substitution axis of economic reasoning has no purchase, while the allocation axis retains all of its.

\subsection{Layered phenomena and distributive reduction}
\label{p02:sec:pe-layered}

The two axes do not partition intimate life into separate provinces; they run through single acts. He brings her the flowers she once mentioned loving. The money that bought them is allocative: transferable, budgetable, its source a matter of indifference to the act's meaning. The hour spent fetching them is allocative in part: a courier could carry the flowers without destroying the act. The remembering of which flowers, the noticing of the moment at which they would matter, is constitutive: not transferable to a courier of cognition without the act's becoming a different act. Intimate acts are, in general, \term{layered phenomena}: an allocative face, where scarcity, transferability, and measurement have full grip, laminated to a constitutive face, where the question is never how cheaply but only by whom. The layering, and not any property of machines, is the deep structure of the entire problem.

Each face has its proper science, and each science errs at the lamination line in its characteristic direction. On the allocative face, allocative reasoning is not merely permitted but ethically required: the predecessor paper's defense of precise household accounting, that the ledger is the rite which prevents the silent asymmetric accumulation of burden, is an allocative-justice argument, and a couple that refused all distributive grammar in the name of love would be committing the inverse error, mystifying real inequities of time and toil that only measurement can surface. The distributive critique of housework, the domestic labour debate of materialist feminism, operates legitimately and importantly on this face.\footnote{See \textcite{delphy1984close} for the materialist analysis; the debate's lesson for this paper is precisely that the allocative face of domestic life is real and that refusing to count is itself a politics.} On the constitutive face, the same grammar misfires, and the misfiring has a distinguished literature of its own: \textcite{hochschild1983managed} documented, in the regime of commercial emotional labor, what becomes of constitutive feeling when it is processed as an allocable input, which makes her the empirical recorder of the category error this section is defining. Between the two errs the largest formal target: \textcite{becker1981treatise} extends the allocative grammar to marriage, altruism, and the family with complete consistency and real success on the allocative face, and with a systematic blindness to the lamination, every constitutive item entering the model, if at all, as a preference parameter, which is to say, as already reduced. The position this paper requires is the one \textcite{zelizer2005purchase} marked out against both ``hostile worlds'' (economy and intimacy as mutually corrupting spheres to be quarantined) and ``nothing-but'' reduction (intimacy as merely exchange): intimate life is shot through with economic practice, and everything depends on the \emph{matching} of practice to relation. The layered-phenomenon analysis is offered as a criterial sharpening of that middle path: the matching norm is that allocative grammar's jurisdiction ends at the lamination line within each act.

The transgression of that line now has a name and a definition. \term{Distributive reduction}: the treatment of a layered act as exhausted by its allocative face; operationally, the optimization of the act's allocative face under which its constitutive face is silently substituted. This is the political-economic restatement of everything \S\ref{p02:sec:synthesis}'s criteria were built to detect, and it is the exact, non-metaphorical content of the Marxian diagnosis imported in \S\ref{p02:sec:marx}: alienation, on the 1844 grammar, \emph{is} the degradation of constitutive activity into allocable means, praxis processed as poiesis, and the AI-mediated case is that degradation executed not by a foreman but by an optimizer.

\subsection{Three operations of the machine upon the budget}
\label{p02:sec:pe-operations}

With the layering in hand, the machine's intervention in the cognitive budget resolves into three distinct operations that the undifferentiated word ``mediation'' had been concealing.

\subsubsection{Reallocation: the liberation effect}
\label{p02:sec:pe-reallocation}

The first operation is the takeover of allocative cognition: ledgers, schedules, logistics, the factual memory of dates and lists. Here the machine spends nothing constitutive; it absorbs budget lines whose expender was always a matter of indifference, and thereby \emph{releases} attention toward the relationship. This is the classical logic of the division of labor applied at the lamination line: the machine performs the poiesis-shaped cognition, the human retains the praxis-shaped cognition, and total constitutive expenditure can rise. The predecessor paper's household modules are hereby vindicated in the opponent's own grammar: their justification, ``return attention to the relation,'' is a Beckerian household-productivity argument \parencite{becker1965theory,becker1981treatise}, and it is sound, because every line it reallocates lies on the allocative face. The liberation effect is real, legitimate, and the strongest truth in the technology's favor; the synthesis must protect it, and does.

\subsubsection{Substitution: the reduction executed}
\label{p02:sec:pe-substitution}

The second operation presents, on the budget sheet, as indistinguishable from the first: a cost reduction on a cognitive line item. The ambient tip lowers the cost of ``attending to her.'' But that line is not like the others, by the boundary theorem of \S\ref{p02:sec:pe-budget}: it is the item whose value to its end is constituted by its cost's being personally borne. Cost-minimizing it is therefore not efficiency but self-defeat, the act of buying, at a discount, an asymptotically worthless good, since at the limit of zero personal expenditure the item's yield, \emph{that he attended}, is also zero. Distributive reduction is this operation: the second masquerading as the first, substitution wearing reallocation's accounting. The four criteria of \S\ref{p02:sec:synthesis-criteria} can accordingly be redescribed, without remainder, as an audit procedure for telling the two operations apart: object, provenance, and initiative locate the line item's face; attribution expectations record whether the books, as presented to the relationship's other member, have concealed the substitution.

\subsubsection{Managerial migration: the Taylorism of love}
\label{p02:sec:pe-management}

The third operation is the subtlest, and it is the one the push-architecture adds. An ambient engine that initiates, ``notice her now, and in this register,'' has taken over not a budget line but the \emph{management of the budget}: the deciding of when, toward whom, and toward which aspect attention shall be spent. This is, structurally, the move \textcite{braverman1974labor} identified as the core of Taylorist deskilling: the separation of conception from execution, planning withdrawn from the worker into the office, the worker retained as the executor of issued instructions. Its appearance here may be called the \term{Taylorism of love}: the engine becomes the planning office of intimate attention, the lover its scheduled operative. The third criterion (direction of initiative, \S\ref{p02:sec:synthesis-criteria}) hereby receives its political-economic name, and the fragment-on-machines argument (\S\ref{p02:sec:marx-fragment}) its mechanism: conception, once migrated, does not exercise the agent's own faculty of noticing, and the faculty, unexercised, atrophies, which is the deskilling that the dynamic criterion (\S\ref{p02:sec:synthesis-dynamic}) watches for. One empirical caution completes the operation's analysis: the liberation effect of \S\ref{p02:sec:pe-reallocation} cannot be presumed even for licit reallocations, since, as \textcite{cowan1983more} showed for household technology, mechanization historically raised the standards of expected output and reabsorbed the liberated time; a tip system may likewise inflate the expected standard of thoughtfulness and conscript the couple into a denser production of gestures. Whether liberation or reabsorption occurs is a question of fact about each household's trajectory, one more matter remanded to the dynamic criterion's standing tribunal.

\subsection{The deflation of the gesture's value-form}
\label{p02:sec:pe-deflation}

The analysis so far has been micro; the deepest yield of the allocative grammar is systemic. Why does the gesture ``he remembered the flowers'' carry thick meaning at all? Because, under given social and technical conditions, producing such a gesture \emph{socially necessarily} requires a real expenditure of personal attention, so that the gesture functions as credible evidence of the expenditure. The structure is precisely that of a costly signal, whose information content is secured by, and only by, its cost \parencite{spence1973job}; and it is equally the logic of the gift, whose value contains the giver's expenditure and is non-delegable for that reason \parencite{mauss1925gift}. Transplant, now, the value-form analysis: as the value of a commodity falls with the socially necessary labor time of its production \parencite{marx1867capital}, so the evidentiary value of a class of gestures falls with the socially necessary \emph{attention} time of their production. When generative mediation collapses the marginal cost of thoughtful gestures toward zero, the currency deflates: the gesture-form ceases to be readable as proof of personal attending, for everyone, including the lovers who never used the machine.

Three consequences follow, each of which the paper has needed. First, the fourth criterion's attribution baseline (\S\ref{p02:sec:synthesis-criteria}) acquires its missing dynamics: the baseline of expected mediation is not free-floating convention but moves with the socially necessary attention time of gesture production, which is why technologies, and only technologies of sufficient diffusion, move it. Second, the historical difference from Cyrano is explained rather than gestured at: the handwritten ballade had so high a social cost that ghostly authorship was rare, individual, and scandalous; zero-marginal-cost mediation converts the same transgression from an episode into a currency regime, which is the precise and non-alarmist content of the claim that AI changes the problem's scale rather than its kind. Third, and dialectically most important: this entire argument has run inside allocative grammar, prices, costs, signals, value-forms, and has arrived, independently, at the constitutive analysis's verdict, that the cheapened gesture is a different and lesser thing. The two grammars, assigned separate jurisdictions in \S\ref{p02:sec:thesis-jurisdiction}, converge at the systemic level, and the convergence joins that of \S\ref{p02:sec:buddhism-convergence} as the strongest species of evidence the paper's method admits.

\subsection{What political economy can and cannot say}
\label{p02:sec:pe-verdict}

The section's result may be stated as a jurisdictional finding, the third application of the method that has organized the paper (first among ethical traditions, then between law and ethics, now between the two faces of a single act). Political economy, granted its full voice, describes the whole outline of the danger: it distinguishes reallocation from substitution, identifies the managerial migration, predicts the deskilling, warns of reabsorption, and derives the systemic deflation that drives the attribution baseline. It cannot say, in its own vocabulary, why the substitution is a \emph{wrong} rather than an arbitrage, for the wrongness lives on the face of the act where its categories do not reach, in the constitution of being-loved by the lover's own attending, and for that the paper's account remains the one given at \S\ref{p02:sec:synthesis-twowrongs} and \S\ref{p02:sec:implications-answer}. The division is not a defeat for either grammar. It is the paper's thesis exemplified: each science whole within its jurisdiction, neither sovereign over the act entire, and the lamination line between them patrolled, as everything here is finally patrolled, by the standing tribunal of practice. One question, however, political economy has raised and not yet been permitted to finish: if the danger can be stated entirely in the vocabulary of costs and substitutions, perhaps it is, after all, only a creature of the social form in which such systems are commercially deployed, and would vanish where that form is absent. The next section conducts the experiment.

% =====================================================================
\section{Exploitation and Ownership: The Eliminative Test}
\label{p02:sec:exploitation}
% =====================================================================

\subsection{What exploitation requires, and the commercial case}
\label{p02:sec:ex-structure}

The classical structure of exploitation has three elements: a surplus is produced; the surplus is appropriated, uncompensated, by another; and the appropriation rests on a structural position of dominance, paradigmatically the monopoly of the means of production. Commercially deployed intimacy systems satisfy the structure without strain. On \textcite{zuboff2019age}'s analysis, the relationship's interior data is rendered as free raw material and the behavioral surplus extracted from it is appropriated by the platform; on the digital-labour analysis \parencite{fuchs2014digital}, the users' engagement is itself unwaged producing activity. None of this is contested here, and none of it is the present subject, by the scope stipulation of \S\ref{p02:sec:intro-method}. What that stipulation now makes possible is a controlled experiment that the commercial literature, for all its force, cannot run: remove capital from the scenario entirely, the system self-built, self-hosted, serving no third interest, and observe what, if anything, remains of the wrong. If the residue is zero, the entire problem of this paper was a problem of capitalism wearing a technological mask, and AI mediation as such stands acquitted. If the residue is not zero, then there exists a stratum of the wrong that exploitation theory cannot capture, and the paper's object is, precisely, that stratum.\footnote{The stipulation may now also be given its political-economic name: a self-built, self-hosted system is not an engineering preference but the deliberate dismantling of the third element, the structural dominance of an appropriating owner; it is the refusal of the platform form. That the predecessor paper's system is of this kind is what qualifies it as the experiment's apparatus.} Two configurations of the capital-free case must be distinguished, and they behave differently.

\subsection{Configuration one: the single builder}
\label{p02:sec:ex-single}

Let one partner build the system, for the relationship, alone. Classical exploitation fails immediately: there is no appropriator; the builder's labor is his own, its product remains inside the relation that generated it; in the vocabulary of generative justice, the value circulates unalienated among its generators. Two candidate residues must nevertheless be examined honestly. The first is self-exploitation in the sense of \textcite{han2015burnout}: the achievement-subject who, no longer commanded by any master, presses himself in the master's place, and who may now extend the regime of self-optimization into love itself, the relationship one more project to be performed. The paper registers this as a real reflexive hazard, particularly for a builder of elaborate relational infrastructure, while noting that its wrong-structure is not appropriative but disciplinary, the internalization of a norm of performance, and that its proper theoretical home is accordingly not exploitation theory; it is held at the margin as a standing caution rather than developed as an axis.

The second residue is the section's finding. Even with no surplus transferred, the single-builder configuration institutes a structural asymmetry: the builder holds what may be called the \term{conception-rights} over the mediating layer, the deciding of how the system prompts, when it is tender, in what register it speaks, which aspects of the relationship it models at all, while the other partner, even fully informed and gladly consenting, occupies the position of the mediated: she lives within a normative apparatus she did not conceive. This is Braverman's separation of conception from execution (\S\ref{p02:sec:pe-management}) rotated ninety degrees: not capital monopolizing conception against labor, but one member of the dyad monopolizing conception of the relation's mediating layer against the other. And it is, precisely, the configuration whose possibility \textcite{roemer1982general} established when he decoupled exploitation from domination and showed that each can occur without the other: what the single-builder case exhibits is \emph{domination without exploitation}, dominance constituted not by the appropriation of surplus but by the unilateral holding of the means of mediation. No value is taken; a position is. The mechanism of this position, how a structure built from love can function as power, exceeds political-economic vocabulary and is remanded to \S\ref{p02:sec:power}.

\subsection{Configuration two: the joint builders}
\label{p02:sec:ex-joint}

Let both partners conceive, design, and own the system together. The configuration is doubly clean: no external appropriator, and the conception asymmetry of \S\ref{p02:sec:ex-single} dissolved, both holding the pencil at the drawing of every default. It deserves to be noticed what this configuration is, in the classical vocabulary: common ownership of the means of (cognitive) production, joint planning of the producing process, the value of the labor circulating entirely among its generators; that is, the form Marx reserved for the freely associated producers, realized at the smallest social unit that can exhibit it.\footnote{One may call it, with strict accuracy and at the risk of levity, a communism of two. The phrase is confined to this footnote.} Moreover the building is not merely the precondition of a relational good but is one: the joint authoring of the mediating layer is itself an exercise of joint attention, a co-creation, a stretch of the relationship's praxis, so that in this configuration the labor of making the instrument belongs to the constitutive face of the relation rather than standing outside it. Alienation, across every register this paper has used, reaches its minimum here.

Its minimum, and not zero. For consider the co-built engine on the evening it performs: the noticing that she has lately gone unappreciated is conceived, tonight, by the apparatus, and the gesture executed by him; and the analysis of \S\S\ref{p02:sec:pe-substitution}, \ref{p02:sec:synthesis-criteria} applies with unreduced force, since nothing in it ever turned on who owned the machine. Joint ownership of the mirror does not make the face in it one's own. The constitutive substitution, where it occurs, occurs identically under platform capitalism, under single building, and under the most symmetrical co-construction that love can design.

\subsection{The eliminative finding}
\label{p02:sec:ex-finding}

The experiment's result can be tabulated.

\begin{center}
\small
\begin{tabular}{lccc}
\hline
Configuration & Exploitation & Domination & Substituted attending \\
 & (appropriated surplus) & (conception asymmetry) & (constitutive face) \\
\hline
Commercial platform & present & present & present \\
Single builder & absent & residual & present \\
Joint builders & absent & absent & \textbf{present} \\
\hline
\end{tabular}
\end{center}

\vspace{0.3em}
\noindent Read column-wise, the table assigns each critical theory its exact jurisdiction: exploitation theory explains the first column's difference between the rows, and the difference is large, real, and reason enough to prefer the lower rows; the theory of domination explains the second column, and will be developed in \S\ref{p02:sec:power}; and the third column does not vary. There exists a wrong that is \term{insensitive to ownership configuration}: it is not exploitation, not appropriation, and not domination, since it survives the elimination of all three; it is the act-philosophical wrong, the ceding of a strictly personal act, exactly as \S\ref{p02:sec:synthesis} characterized it, and the experiment has now shown by elimination that no political-economic reform, not even the perfect one, reaches it. This is the strongest form of the paper's central claim, and it could only have been obtained by granting the political-economic tribunals everything they asked.

\subsection{The generative form: from elimination to program}
\label{p02:sec:ex-generative}

The eliminative test states which reforms do not reach the wrong. It does not state what right practice looks like, and a paper whose diagnostic registers are now five deep owes its reader the constructive complement. The complement exists, and the predecessor framework already stands in working contact with it: the theory of \term{generative justice} \parencite{eglash2016introduction}, whose critique of distributive justice is structurally kin to this paper's critique of imputation (\S\ref{p02:sec:thesis-jurisdiction}). Distributive justice, on Eglash's analysis, is a top-down, ex post correction of flows already alienated at the point of extraction; generative justice asks instead that value circulate \emph{unalienated} among those who generate it, under their own governance, from the start. Of the three value-forms the theory tracks, unalienated labor value, ecological value, and \emph{expressive} value, the third is decisive here, for the currency in which expressive value circulates, in the open-source communities that are the theory's paradigm, is precisely \emph{attribution}: the fourth criterion of \S\ref{p02:sec:synthesis-criteria} had, it turns out, a theory of justice waiting for it.

Applied to the dyad, the question ``how is value returned to its generators'' resolves into four layers, and the resolution discloses that the remedies this paper has scattered across its sections are not a list but a single practice. At the \emph{data} layer, the generator of the relational sediment is the relation itself, and the practice is the one \S\ref{p02:sec:ex-structure} recorded: self-building and self-hosting are the refusal to let the sediment leave the circle of its generation; the scope stipulation of this paper may now be given its third name, after engineering choice and political-economic refusal: it is generative justice at the layer of data. At the \emph{attribution} layer the analysis yields a point not yet stated in any prior register: ask who generates the value of the gesture, and the answer includes \emph{her}, for the engine's inference runs on her moods, her offhand remarks, her self-disclosure, so that the recipient of the tender tip is the \term{uncredited co-generator} of the gesture directed at her; her generated value is extracted, processed, and returned to her under another's authorship. The supply chain of thoughtfulness is thereby restated in generative vocabulary: the wrong is not that value leaves the relation, it circulates back, but that it circulates with its authorship overwritten, expressive value alienated not in destination but in name. The remedy, attribution hygiene, the alignment of the credit flow with the generation flow, gives the fourth criterion and the transparency cure of \S\ref{p02:sec:synthesis-twowrongs} a positive justice-theoretic ground, no longer merely the avoidance of deception but the restoration of expressive value to its sources. At the \emph{capacity} layer, the theory's signature concept, recursion, the property by which a generative cycle deepens the generative capacity of its own participants \parencite{eglash1999african,eglash2014recursive,eglash2019ai}, soil that grows richer, contributors that grow abler, furnishes the dynamic criterion of \S\ref{p02:sec:synthesis-dynamic} with its fourth independent derivation, after the Hegelian return, the Buddhist raft, and the Bravermanian mechanism: a tip that trains the lover's own noticing is a generative cycle, capacity returning enriched to its generator; a tip on which reliance deepens while the native faculty recedes is a broken cycle, capacity accumulating, alienated, on the apparatus side. And at the \emph{constitutional} layer, generative justice's insistence on self-governance, that those regulated by a value structure participate in generating it, is exactly the co-constitution that \S\ref{p02:sec:power} will demand: the single-builder configuration alienates the one value the theory treats as sovereign, the generation of the structure itself. Four layers, one principle: the data back to the relation, the name back to the generators, the capacity back to the agent, the constituent power back to both.

One objection must be met, because it would otherwise let the generative register quietly annex the constitutive face. A purist of circulation might argue: correct the credit flows, let the gesture arrive labeled with its true supply chain, machine and her own sediment included, and the circulation is unalienated, the wrong zero. The reply is the paper's core, now in generative vocabulary: the value at issue in the constitutive face, \emph{that it was he who attended}, cannot be generated by any other node, however honestly credited, for generative justice presupposes identifiable generators of circulable value, and the analysis of strictly personal acts marks exactly the class of value that admits no proxy generation; where the substitution occurs, the advertised good is not misrouted but nonexistent (\S\ref{p02:sec:implications-answer}). Generative justice therefore takes its place in the architecture not as a fourth column of the table above but as the program that turns the first three columns green; the fourth column remains guarded by the dynamic criterion, which is, in the theory's own vocabulary, nothing other than its recursion principle standing watch.

% =====================================================================
\section{Power: The Architectural Constitutionalization of Intimacy}
\label{p02:sec:power}
% =====================================================================

\subsection{The third face of power in the home}
\label{p02:sec:power-faces}

The residual of the single-builder configuration, domination without exploitation, now receives its mechanism, and the mechanism is what makes it dangerous: it is implicit. Explicit asymmetries of power in intimate life, who earns, who decides the large purchases, who concedes in open conflict, are visible, contestable, and negotiable; their visibility is the door through which renegotiation enters. The conception asymmetry has a different mode of existence. The builder's values do not appear in the relationship as claims, which could be answered, but as \emph{defaults}: as the timing of prompts, the warmth of a register, the boundaries of the modules, the very inventory of what about the relationship is modeled and what is not. The other partner does not encounter a position she might dispute; she inhabits an already-constituted space. Three features compound the implicitness. The structure operates in the form of the good: it is the product not of a will to control but of love and care, and benevolence is its cloak of invisibility, since to contest a tender default is to appear to contest the tenderness, at an emotional price that forecloses the contest before it forms. It is not experienced as power: what the mediated partner phenomenally encounters is that the system is thoughtful, not that she lives inside another's legislated norm-space. And it shapes the agenda and the preferences themselves: in \textcite{lukes2021power}'s terms this is power in its third dimension, beyond open decision (the first) and beyond agenda control \parencite{bachrach1962two}, though it exercises that too, the module boundaries determining which relational questions can so much as surface; it is the slow formation of both partners' sense of what good intimacy is, conducted by an artifact one of them authored.

\subsection{Code as the law of the home}
\label{p02:sec:power-code}

The theoretical frame for such a structure exists, built for the public sphere and never, to the author's knowledge, transposed to this one. \textcite{lessig2006code}: in digitally mediated spaces, architecture regulates alongside law, market, and norms, and regulates without needing to be perceived as regulation. \textcite{winner1980artifacts}: technical arrangements are themselves political arrangements, the bridge too low for the bus is a policy with no author on record. Transpose both into the intimate domain and the single-builder system shows itself as what it is: \term{the architectural constitutionalization of intimacy}, a constitution for the relationship, drafted by one party, conferring no amendment rights on the other, and not presenting itself as a constitution at all. The predecessor paper's standing design-philosophy documents, each clause of which is individually admirable, are, in this light, the material record of a unilateral constituent act: one person's legislation of how we are to be mediated. The transposition also yields a corollary about the constitutive unsaid that the predecessor paper placed among intimacy's defining features: the relational-contract insight \parencite{macneil1980new} was that ongoing relations are governed by implicit norms that the parties evolve together and never fully articulate; an architected mediating layer converts a tract of those living, jointly drifting norms into fixed, unilaterally authored, executable text. This is a new species of injury to the unsaid: not the saying of it, which the predecessor paper guarded against, but the \emph{writing of it for the other}, without her hand.

\subsection{The weakness of subsequent consent, and constitution as a verb}
\label{p02:sec:power-consent}

It will be objected that disclosure and consent dissolve the difficulty: she learns of the system, approves it, uses it gladly. The objection fails for a reason older than the technology. Hume's argument against the original contract \parencite{hume1748original} was that consent given after the institutions are built, by one who must live somewhere, is not the authorship of those institutions: the sailor who wakes aboard does not, by remaining, become the ship's designer. Subsequent consent to an architected intimacy is habitation, not constituent power; the constitutional moment has passed, and it passed with one signature on it. This, and not only the co-creative good of building together, is what gives the joint configuration of \S\ref{p02:sec:ex-joint} its full normative weight: it is the only configuration in which constituent power is symmetric, and it converts constitution from a noun, a finished thing one party delivers, into a verb, the continuous joint constituting of the shared mediating layer, with amendment rights that never lapse. Two consequences must be drawn without flinching. First, reflexively: the touchstone system of this paper and its predecessor was, as a matter of history, single-built; the analysis of this section therefore indicts, in measured part, the very apparatus from whose examination the paper arose, and the honesty norms governing both papers require that this be said in the text rather than survived in silence; the indicated remedy is not demolition but co-constitution, the retroactive and ongoing extension of the drafting hand. Second, doctrinally: the praise extended in \S\ref{p02:sec:thesis-jurisdiction} to the restraint traditions of family law must now be qualified, for restraint toward the interior of the home is also non-patrol of the home's third-dimensional power, and the architectural constitution is exactly the kind of structure that judicial restraint cannot see. The classic feminist correction, that the personal is political, here acquires a precise technical content, and with it the analysis passes to the tribunal that has been owed the floor longest.

% =====================================================================
\section{The Feminist Hearing}
\label{p02:sec:feminism}
% =====================================================================

\subsection{Debts, and a repayment}
\label{p02:sec:fem-debts}

This paper has been borrowing from feminist theory since its criteria were first stated: the placement of attentiveness as care's first moment is Tronto's; the empirical record of constitutive feeling processed as allocable input is Hochschild's; the insistence that the home's allocative face be counted is the materialist feminism of the domestic labour debate. The audit is owed, and the first repayment can be made at once, because the paper's lamination-line analysis (\S\ref{p02:sec:pe-layered}) is a late contribution to a controversy feminism conducted internally half a century ago. The wages-for-housework campaign \parencite{federici1975wages} demanded that the labor of the home be counted, waged, made visible as labor; its critics, inside the movement, feared that pricing care would commensurate and so corrupt it; feminist economics subsequently built the accounting that makes invisible care countable \parencite{waring1988if,folbre2001invisible}. The controversy was, in this paper's vocabulary, a dispute over where the lamination line runs: the campaign was right that the allocative face of domestic labor is real, that refusing to count it is itself a politics, and that the refusal had a beneficiary; the critics were right that the constitutive face does not survive commensuration. The layered analysis is offered back to that debate as a criterial settlement: count everything on the allocative face, and refuse the reduction of the constitutive face, in the same household and without contradiction.

\subsection{The gendered budget: mental load, redistribution, and the counterfeit answer}
\label{p02:sec:fem-load}

The cognitive budget of \S\ref{p02:sec:poleconomy} is not, in fact, gender-neutral, and the empirical literature locates the skew at the exact joint this paper has been examining. \textcite{daminger2019cognitive}, decomposing household cognitive labor into anticipation, identification, decision, and monitoring, finds the anticipatory and monitoring stages, that is, precisely the \emph{inception of noticing}, the conception-stage of domestic and relational attending, disproportionately performed by women; the mental load is, in the paper's terms, a gendered distribution of exactly the budget-management function whose migration \S\ref{p02:sec:pe-management} analyzed. The tender tip therefore does not enter a neutral field. It enters a field already structured by a standing redistributive claim, decades old: that he should carry his share of the noticing. And here the analysis yields what may be its sharpest single result. The feminist demand is for \emph{redistribution} of attending, from her to him; what the machine offers is \emph{substitution} of attending, from him to it. The two are indistinguishable on the allocative face, her measured load falls either way, and categorically different on the constitutive face: he finally noticed, except that it was not he who noticed. An equality claim is answered in counterfeit currency; the distributive reduction of \S\ref{p02:sec:pe-layered} acquires its gendered form, and the four criteria acquire a further office, that of distinguishing genuine redistribution (he takes up the noticing, perhaps scaffolded at first, the scaffold falling away on the dynamic criterion's good branch) from automated simulation of it. Honesty requires the converse to be stated with equal force: where the mental load is crushing and redistribution is in fact unavailable, the automation of the load's genuinely allocative items, the lists, the logistics, the calendar of obligations, is real relief and not counterfeit anything, which is the liberation effect of \S\ref{p02:sec:pe-reallocation} doing exactly what it should; the wrong is specific, the substitution of the constitutive inception while the books show redistribution, and the criteria exist to find it.

\subsection{Who legislates}
\label{p02:sec:fem-legislates}

The constituent power of \S\ref{p02:sec:power} is likewise not distributed by coin-flip. The capacity to build, configure, and maintain the mediating layer tracks the distribution of technical capital, which is, at present and in fact, gendered; the architect of the home's constitution will, in the statistical run of cases, be the male partner, and the architectural constitutionalization of intimacy will land, in the statistical run of cases, as one more layer on an already tilted field, its defaults encoding one party's conception of care in the medium least available to the other's amendment. This is a claim about tendencies and not essences, and it sharpens rather than alters the section's prescriptions: the case for co-constitution (\S\ref{p02:sec:power-consent}) is also, in the actual world, a case about gender; and the reflexive admission recorded there is to be read accordingly, since the author's position in this analysis is the statistically typical one. The correction this tribunal administers to \S\ref{p02:sec:thesis-jurisdiction} has already been entered; it bears repeating in the tribunal's own voice: a jurisprudence that prides itself on not entering the home should know what it is thereby declining to patrol.

\subsection{A positive theory of the constitutive face}
\label{p02:sec:fem-positive}

Every defense of the constitutive face so far mounted has been negative: freedom from alienation, from Taylorization, from deflation, from unilateral constitution. The tradition that has always held this terrain supplies what the paper still lacks, an affirmative account. In \textcite{noddings1984caring}, the root of the caring relation is engrossment: a receptive, non-selective attention to the concrete other, displacement of one's own frame into hers; and this attention is not an instrument by which caring is produced, it \emph{is} the caring, the relation's mode of existence and not its means. \textcite{kittay1999love} grounds the same point in the labor of dependency: the work of attending to a particular vulnerable other is not a fungible service that happens, regrettably, to be hard to delegate; its non-fungibility is internal to what it is. Herein the eliminative residue of \S\ref{p02:sec:ex-finding}, the wrong that survived every ownership configuration, receives at last its positive name, and the paper's longest arc closes: the constitutive face needed no exemption from economics, no indulgence from law, and no nostalgia from anyone, because it has its own jurisprudence, the ethics of care, in which the non-delegability of attending is not a restriction on love's efficiency but the definition of love's presence. The feminist hearing thus performs double duty in the architecture of the whole: it audits every prior tribunal, the political economy for the gender of its budget, the power analysis for the gender of its constituent, the law for the politics of its restraint, and it supplies the affirmative ground on which the synthesis of \S\ref{p02:sec:synthesis}, drafted in doctrinal idiom, was silently standing all along.

% =====================================================================
\section{Implications}
\label{p02:sec:implications}
% =====================================================================

\subsection{The answer returned to the predecessor paper}
\label{p02:sec:implications-answer}

The predecessor paper asked whether a genuinely improved relation can redeem an alienated process, and kept the question open. The answer this paper returns is: no, and the question was subtly ill-posed, and both halves matter. No: because improvement and alienation are not on one scale, redemption is a transaction between commensurables, and the process defect, where the grounds and the dynamic criterion find it, is not a cost that outcomes offset but a change in what the outcome \emph{is}. In intimate life, process partially constitutes outcome: the content of ``being loved'' includes ``that it is he who attends.'' Where the constitutive links are substituted, the outcome ``the relationship improved'' has not been purchased at a regrettable price; it has been quietly redefined, the dyad become a triad, and the improvement, though real, is an improvement of the redefined thing, which is why the stipulated unity of subjective and objective betterment never settled the matter. In the vocabulary that \S\ref{p02:sec:pe-layered} made available: the reports of improvement issue from the act's allocative face, where everything has indeed become cheaper and smoother, while the substitution occurred on the constitutive face, which no report drawn from the first face can audit; the open question was, in its deep structure, an invitation to consummate a distributive reduction, and the correct response to such an invitation is not acceptance or refusal but the exposure of the reduction. The consequentialist defense does not lose the argument; its object of evaluation, an outcome independent of process, fails to exist in this domain. The exposure of this ill-posedness is not an evasion of the original question but its resolution: what the question needed was not an answer but a diagnosis, and the dialectical method's verdict on its own opening move is the most Hegelian result in the paper.

And yet the predecessor paper's provisional tolerance of its tender-tip module, under the four constraints, is largely vindicated rather than revoked, for the constraints turn out to have been criteria in embryo: opt-in and infrequency preserved initiative (criterion three); occasion-and-register-only confined the object of mediation (criterion one); the prohibition on presenting outputs as knowledge of her mind policed attribution (criterion four); refusability and non-recording kept the raft loosely held (the dynamic criterion). What this paper adds is the theory of why those particular constraints, and the warning the theory carries: the constraints hold the module at the boundary of the exception domain, not outside it, and boundaries, under the dynamic criterion, are not places to live but places to keep watch.

\subsection{The boundary of consent in intimate life}
\label{p02:sec:implications-consent}

The two-wrongs architecture yields a determinate position on the unwaivable core. Intimate life does contain one, but it is narrower than the romantic and broader than the liberal intuition: it comprises not gestures, contents, or technologies, but the constitutive links of intimate attending, the inception of noticing, the interpretation of the particular other, as exercised capacities of the persons in relation. Partners may consent to any amount of mediation of means; their consent fully cures the attributive wrong; what their agreement cannot do is make the substituted inception of attending be theirs, any more than the testator's consent can make another's will his testament. This is not paternalism toward lovers; it is the same recognition the law has always extended to testators and spouses: that some acts are constituted by their personal performance, and that respecting persons includes refusing the fiction that they can be elsewhere while performing them. At the scale of social baselines, the same architecture issues a caution rather than a prohibition: ambient intimacy-mediation at zero cost will press the attribution baseline toward including it, as technologies always have and by the deflationary mechanism identified at \S\ref{p02:sec:pe-deflation}, and the baseline's movement will cure the attributive wrong wholesale; but no movement of expectation can touch the grounds, and a society whose baseline has normalized substituted noticing has not abolished the process wrong, only its visibility. The paper registers this as the precise, bounded sense in which something true survives from the old romantic protest against the machine.

\subsection{The dialectical reservation}
\label{p02:sec:implications-reservation}

Fidelity to the method requires the synthesis to state its own non-finality. The grounds are graded; the baseline moves; the dynamic criterion is legible only across time and only in particulars. The exception domain here constructed is therefore not a fence but a jurisdiction: a standing assignment of the question, site by site, gesture by gesture, year by year, to the tribunal the predecessor paper called dialectical-positivist practice, which measures, describes, and asks again, and never finally adjudicates. To its docket the later sections have added a second standing matter: the constitution of the mediating layer itself, which \S\ref{p02:sec:power-consent} showed can never be ratified once and inhabited thereafter, but must remain jointly amendable, constitution as a verb, the drafting hand permanently shared. The docket may equally be read in generative vocabulary: what the tribunal audits, season by season, is whether the relation's cycles still deepen their generators, the data, the name, the capacity, and the constituent power each returning to its source. The paper's contribution is to have given that tribunal a doctrine to apply, four criteria, two layers, one dynamic test, two wrongs, two remedies, a constitutional question never closed, and a generative program for the whole, not to have relieved it of session.

% =====================================================================
\section{Limitations}
\label{p02:sec:limits}
% =====================================================================

Five limitations are stated without attenuation. First, the empirical premises of the dynamic criterion, deskilling of attentiveness under ambient suggestion, capacity migration under pulled use, are here motivated theoretically (\S\ref{p02:sec:marx-fragment}, \S\ref{p02:sec:buddhism-raft}) and by analogy to the AI-mediated-communication literature on perception and trust \parencite{hancock2020ai,jakesch2019ai}, but they are not yet established for intimate dyads; the same conditionality attaches to the two political-economic mechanisms of \S\ref{p02:sec:poleconomy}, the deflation of the gesture's value-form and the Cowan-style reabsorption of liberated attention, which are theoretically derived mechanism hypotheses and empirically open; the paper's verdicts on the failing branch are conditional on facts that longitudinal study must supply. Second, the comparative legal survey of strictly personal acts is in this draft a sketch resting on a small number of verified anchors; the six-jurisdiction examination the synthesis deserves is deferred to the next revision, and all statutory citations carry verification flags until then. Third, the Marxian apparatus is used in its philosophical-anthropological lineage by declared choice; readers who hold that alienation is meaningless outside the critique of political economy will find the antithesis re-described rather than refuted, and the paper has argued for, but cannot compel, the legitimacy of the structural transplant. Fourth, the Buddhist materials are drawn from the P\=ali corpus in translation and deployed philosophically, not philologically; the paper claims convergence of structure, not exegetical completeness, and parallel examinations from the Confucian and Daoist traditions, each of which bears directly on the question and was deliberately reserved, remain future work. Fifth, the analysis is built on the dyadic intimate relationship as characterized in the predecessor paper; its extension to filial, parental, and friendship intimacies, where role-definition and asymmetry alter the grounds, is unexamined here and may not be conservative. Sixth, the three later tribunals are convened at sketch depth relative to the first three: the eliminative test reasons over idealized configurations and abstracts from mixed and migrating ownership forms; the power analysis transposes frameworks built for public architectures and owes a fuller account of resistance, exit, and renegotiation inside intimacy; and the feminist hearing rests its empirical premise on a literature, exemplified by \textcite{daminger2019cognitive}, drawn largely from contemporary heterosexual couples in one national context, so that its claims about the gendering of the cognitive budget and of constituent power are claims about documented statistical tendencies within that scope, not about essences, and their generalization across cultures and relationship forms is an open empirical question.

% =====================================================================
\section{Conclusion}
\label{p02:sec:conclusion}
% =====================================================================

Roxane was not wrong that she was loved; she was wrong about where the loving was performed, and the play's grief is that the second error poisoned the first truth. This paper has tried to say, with the precision of six tribunals, dialectical, legal, Buddhist, political-economic, constitutional, and feminist, what kind of error that is and when its modern, ambient, scalable descendant repeats it. Mediation as such is innocent: spirit, love, and law are made of it. The wrong, where there is one, lives in a small and nameable place: in the substitution of the constitutive links of intimate attending, the inception of noticing, the interpretation of the particular other, performed on the relation's own interior sediment, against a baseline that still reads such gestures as personal; and it divides into a wrong of attribution, which honesty cures, and a wrong of process, which nothing cures but the return of the capacity to its agent, the raft released, the training wheels off, the lover restored to his own attending. The eliminative experiment showed this wrong to survive every reform of ownership, including the perfect one; the constitutional analysis showed that even the innocent remainder must be jointly and perpetually authored, the home's code drafted by every hand that must live under it; and the ethics of care supplied the affirmative ground beneath it all, that the attending is not the means of the love but its mode of existence. Between the universal acquittal that the instrumentality doctrine demanded and the universal conviction that the romantic protest desired, the truth has turned out to be doctrinal in the oldest sense: a default rule, an exception domain, criteria, and a tribunal that never adjourns.

% =====================================================================
\section*{Acknowledgments and Disclosure}
% =====================================================================

The author's first acknowledgment is owed to a woman he deeply loves: like a forest, pure, and wise, and kind, and altruistic; to him she is, and will always be, the ``forest girl'' he loves. To have met her is at once miracle and fortune, and the thinking of this paper was also done for her. May all who love find their way to one another, be happy and fortunate, and grow old together. {\cjkfont 祝天下有情人终成眷属,幸福幸运,白头偕老。}

The author discloses that an AI assistant was used in drafting and editing this paper, including the development of formulations under the author's direction across iterative sessions. The conceptions, the dialectical positions taken, and all final judgments are the author's, who bears sole responsibility for the content, including errors. The paper analyzes AI mediation of intimacy while having been drafted with AI mediation; the author considers the practice consistent with the paper's own criteria, scholarly drafting not being a strictly personal act in the sense herein defined, and the disclosure itself being what the fourth criterion requires.
